\section{Mathematical Analysis}
\label{app:theory-full}
\providecommand{\tr}{\operatorname{tr}}
\providecommand{\sym}{\operatorname{sym}}
\providecommand{\skw}{\operatorname{skew}}
\providecommand{\diag}{\operatorname{diag}}
\providecommand{\blkdiag}{\operatorname{blkdiag}}
\providecommand{\St}{\operatorname{St}}
\providecommand{\Hor}{\mathcal H}
\providecommand{\Ver}{\mathcal V}
\providecommand{\EE}{\mathbb E}
\providecommand{\PP}{\mathbb P}
\providecommand{\RR}{\mathbb R}
\providecommand{\op}{\mathrm{op}}
\providecommand{\Id}{\mathrm{Id}}
\providecommand{\sgn}{\operatorname{sign}}
The analysis follows the two interacting parts of STRATA: learning a flag from
gradient statistics and updating parameters with state expressed in that flag.
We first establish the geometry and covariance-to-tracking estimates, then
propagate state errors into descent and sample sensitivity. The final two
subsections treat current-gradient normalization directly and optimize a
directional generalization bound. The linear surrogate, the implemented
adaptive map, and the Gaussian posterior each serve a different role; their
assumptions are stated at the point where each is used.
Three structural reductions sharpen the estimates: prefix overlap becomes a
transportation cost, spectral centering removes isotropic curvature, and
inverse binary KL preserves the empirical-risk dependence. We distinguish
these exact identities and sharp comparison constants from conditional
envelopes for covariance approximation and optimizer-state variation.

\subsection{Flag quotient, tangent decomposition, and retractions}
The geometric variable is a chain of nested subspaces. Rotating a basis inside
one block leaves that chain unchanged, which motivates the quotient and its
horizontal tangent space. The orthogonal block decomposition also separates
captured gradient energy from the residual exactly. Coordinatewise adaptive
moments require additional information about the chosen basis; their behavior
is examined in Proposition~\ref{prop:tf-gauge}.

\begin{definition}[Ambient spaces and block selectors]
\begin{gather}
 \langle X,Y\rangle_F:=\tr(X^\top Y),\quad
 \|X\|_F^2:=\langle X,X\rangle_F,\quad
 \|X\|_{\op}:=\sup_{\|v\|_2=1}\|Xv\|_2,\\
 \sym A:=\tfrac12(A+A^\top),\qquad
 \skw A:=\tfrac12(A-A^\top),\\
 0=r_0<r_1<\cdots<r_K=R<d\le m,\quad
 q_j:=r_j-r_{j-1},\quad q_{K+1}:=d-R,\\
 I_{K+1}:=\{R+1,\ldots,d\},\\
 I_j:=\{r_{j-1}+1,\ldots,r_j\},\quad
 E_j:=\diag(\mathbf1_{I_j})\in\RR^{R\times R},\quad
 F_k:=\sum_{j\le k}E_j,\\
 \St(d,R):=\{U\in\RR^{d\times R}:U^\top U=I_R\},\\
 \mathcal G:=\prod_{j=1}^K\mathrm O(q_j),\qquad
 U\sim V\ \Longleftrightarrow\ 
 V=U\diag(O_1,\ldots,O_K),\quad O_j\in\mathrm O(q_j),\\
 \mathfrak F(q_1,\ldots,q_{K+1})
 :=\St(d,R)/\mathcal G
 \simeq\mathrm O(d)/\prod_{j=1}^{K+1}\mathrm O(q_j),\\
 U_j:=U_{:,I_j},\quad
 P_k(U):=UF_kU^\top,\quad
 \Pi_j(U):=UE_jU^\top,\quad \Pi_{K+1}(U):=I_d-UU^\top.
 \label{eq:tf-projectors}
\end{gather}
\end{definition}
\begin{proposition}[Projector algebra]\label{lem:tf-energy}
\begin{gather}
 P_k^\top=P_k=P_k^2,\quad \tr P_k=r_k,\quad
 P_kP_l=P_{\min(k,l)},\\
 \Pi_i^\top=\Pi_i,\quad \Pi_i\Pi_j=\delta_{ij}\Pi_i,
 \quad\sum_{j=1}^{K+1}\Pi_j=I_d,\\
 G=\sum_{j=1}^{K+1}G\Pi_j,\quad
 \langle G\Pi_i,G\Pi_j\rangle_F=0\quad(i\ne j),\\
 \|G\|_F^2=\sum_{j=1}^K\|GU_j\|_F^2+
 \|G\Pi_{K+1}\|_F^2,\\
 \|G-GP_k\|_F^2=\tr(G^\top G)-\tr(P_kG^\top G).
\end{gather}
\end{proposition}
\begin{proof}
\begin{align}
 P_kP_l&=UF_k(U^\top U)F_lU^\top=UF_{\min(k,l)}U^\top,\\
 \Pi_i\Pi_j&=UE_iE_jU^\top=\delta_{ij}UE_iU^\top,\quad i,j\le K,\\
 (I-UU^\top)UE_j&=0,\qquad
 \sum_{j\le K}UE_jU^\top=UU^\top,\\
 \langle G\Pi_i,G\Pi_j\rangle_F
 &=\tr(\Pi_iG^\top G\Pi_j)
 =\tr(\Pi_j\Pi_iG^\top G)=0,\\
 \|GU_jU_j^\top\|_F^2
 &=\tr(U_jU_j^\top G^\top GU_jU_j^\top)
 =\tr(U_j^\top G^\top GU_j),\\
 \|G(I-P_k)\|_F^2
 &=\tr((I-P_k)G^\top G(I-P_k))
 =\tr(G^\top G)-\tr(P_kG^\top G).
\end{align}
\end{proof}
\begin{proposition}[Vertical and horizontal spaces]\label{prop:rigid-horizontal}
\begin{align}
 T_U\St(d,R)&=\{Z:U^\top Z+Z^\top U=0\},\\
 \Ver_U&=\{U\Omega:\Omega^\top=-\Omega,
                  \ \Omega=\blkdiag(\Omega_{11},\ldots,\Omega_{KK})\},\\
 \Hor_U&=\{Z\in T_U\St(d,R):E_jU^\top ZE_j=0\ (1\le j\le K)\},\\
 T_U\St(d,R)&=\Hor_U\mathbin{\perp}\Ver_U,\\
 \mathsf H_U(A)&=A-U\sym(U^\top A)
   -U\blkdiag\bigl(\skw(U^\top A)\bigr),
 \label{eq:rigid-horizontal}\
 \dim\mathfrak F&=dR-\frac{R^2+\sum_{j=1}^Kq_j^2}{2}.
\end{align}
\end{proposition}
\begin{proof}
\begin{align}
 0&=\left.\frac{\mathrm d}{\mathrm ds}(U+sZ)^\top(U+sZ)\right|_{s=0}
   =U^\top Z+Z^\top U,\\
 Z&=U\Omega+U_\perp B,\quad
 \Omega=U^\top Z=-\Omega^\top,\quad B=U_\perp^\top Z,\\
 \langle Z,U\Lambda\rangle_F
 &=\tr(\Omega^\top\Lambda)
 =\sum_j\tr(\Omega_{jj}^\top\Lambda_{jj}),\\
 Z\perp\Ver_U&\Longleftrightarrow\Omega_{jj}=0\quad(\forall j),\\
 A&=\underbrace{U\sym(U^\top A)}_{(T_U\St)^\perp}
 +\underbrace{U\blkdiag(\skw(U^\top A))}_{\Ver_U}
 +\underbrace{\mathsf H_U(A)}_{\Hor_U},\\
 \langle U S,U\Omega\rangle_F&=\tr(S\Omega)=0
 \quad(S^\top=S,\ \Omega^\top=-\Omega),\\
 \dim\Hor_U&=dR-\frac{R(R+1)}2-\sum_j\frac{q_j(q_j-1)}2.
\end{align}
\end{proof}
We use projector distances to compare subspaces independently of their basis
representatives. The aligned-basis distance is convenient for moment transport;
the following comparisons allow estimates in these two representations to be
combined with explicit constants.

\begin{definition}[Quotient metric and projector distances]
\begin{gather}
 g_{[U]}(\xi,\eta):=\langle\xi_U^{\Hor},\eta_U^{\Hor}\rangle_F,\\
 d_P(U,V)^2:=\sum_{k=1}^K\|P_k(U)-P_k(V)\|_F^2,
 \label{eq:tf-distance}\\
 d_B(U,V)^2:=\sum_{j=1}^K\min_{O_j\in\mathrm O(q_j)}
                    \|U_j-V_jO_j\|_F^2,\\
 d_\Pi(U,V)^2:=\sum_{j=1}^K\|\Pi_j(U)-\Pi_j(V)\|_F^2,\\
 \overline d_\Pi(U,V)^2:=\sum_{j=1}^{K+1}\|\Pi_j(U)-\Pi_j(V)\|_F^2.
\end{gather}
\end{definition}
\begin{lemma}[Overlap transport identity and sharp metric comparisons]\label{lem:rigid-metrics}
For $T_{ja}:=\tr(\Pi_j(U)\Pi_a(V))$, $1\le j,a\le K+1$, set
$W:=\sum_{j\ne a}T_{ja}$. Then
\begin{gather}
 d_P^2=\sum_{j,a}|j-a|T_{ja},\quad\overline d_\Pi^2=2W,\\
 d_B^2\le d_\Pi^2\le2d_B^2,\qquad
 d_\Pi^2\le\overline d_\Pi^2\le2d_P^2,\\
 d_P^2\le\frac K2\overline d_\Pi^2\le Kd_\Pi^2\le2K d_B^2.
\end{gather}
The constants in $\overline d_\Pi^2\le2d_P^2$ and
$d_P^2\le(K/2)\overline d_\Pi^2$ are attained by adjacent-block and
first-to-residual rotations, respectively.
\end{lemma}
\begin{proof}
\begin{align}
 U_j^\top V_j&=L_j\diag(\sigma_{j1},\ldots,\sigma_{jq_j})R_j^\top,
 \quad0\le\sigma_{ji}\le1,\\
 \min_{O_j}\|U_j-V_jO_j\|_F^2&=2\sum_i(1-\sigma_{ji}),\\
 \|\Pi_j(U)-\Pi_j(V)\|_F^2&=2\sum_i(1-\sigma_{ji}^2),\\
 1-\sigma&\le1-\sigma^2\le2(1-\sigma),\quad\sigma\in[0,1],\\
 \sum_aT_{ja}&=q_j,\quad\sum_jT_{ja}=q_a,\\
 \|P_k(U)-P_k(V)\|_F^2
 &=\sum_{j\le k<a}T_{ja}+\sum_{a\le k<j}T_{ja},\\
 \sum_k\|P_k(U)-P_k(V)\|_F^2&=\sum_{j,a}|j-a|T_{ja},\\
 \sum_j\|\Pi_j(U)-\Pi_j(V)\|_F^2&=2\sum_{j\ne a}T_{ja}=2W,\\
 \|\Pi_{K+1}(U)-\Pi_{K+1}(V)\|_F^2
 &=2\sum_{a\le K}T_{K+1,a}\le W.
\end{align}
Thus $W\le d_P^2\le KW$ and $d_\Pi^2\ge W$. A two-plane rotation between
blocks $a<b$ has $W=2\sin^2\theta$ and $d_P^2=(b-a)W$, proving the equality
claims. Keeping the overlap flow avoids repeated triangle inequalities for
the nested prefixes.
\end{proof}
\begin{lemma}[Infinitesimal projector metric]\label{lem:rigid-differential}
For $Q=[U,U_\perp]\in\mathrm O(d)$, $\Omega^\top=-\Omega$,
$\Omega_{jj}=0$, and $U(s)=(Qe^{s\Omega})_{:R}$,
\begin{align}
 \dot P_k(0)&=Q[\Omega,\widehat F_k]Q^\top,\quad
 \widehat F_k:=\diag(I_{r_k},0_{d-r_k}),\\
 \sum_k\|\dot P_k(0)\|_F^2
 &=2\sum_{1\le a<b\le K+1}(b-a)\|\Omega_{ab}\|_F^2,\\
 \|\dot U(0)\|_F^2
 &=2\sum_{a<b\le K}\|\Omega_{ab}\|_F^2
   +\sum_{a\le K}\|\Omega_{a,K+1}\|_F^2,\\
 \|\dot U(0)\|_F^2
 &\le\sum_k\|\dot P_k(0)\|_F^2
 \le2K\|\dot U(0)\|_F^2.
\end{align}
\end{lemma}
\begin{proof}
\begin{align}
 P_k(s)&=Qe^{s\Omega}\widehat F_ke^{-s\Omega}Q^\top,\\
 [\Omega,\widehat F_k]_{ab}
 &=(\mathbf1_{b\le k}-\mathbf1_{a\le k})\Omega_{ab},\\
 \|[\Omega,\widehat F_k]\|_F^2
 &=2\sum_{a\le k<b}\|\Omega_{ab}\|_F^2,\\
 \sum_{k=1}^K\mathbf1_{a\le k<b}&=b-a,\\
 \|(Q\Omega)_{:R}\|_F^2
 &=\|\Omega_{1:K,1:K}\|_F^2+\|\Omega_{K+1,1:K}\|_F^2.
\end{align}
\end{proof}
\begin{proposition}[Polar and QR retractions]\label{prop:rigid-retract}
For $Z\in T_U\St(d,R)$,
\begin{gather}
 X:=U+Z,\quad X^\top X=I+Z^\top Z,\quad\sigma_{\min}(X)\ge1,\\
 R_U^{\mathrm{pol}}(Z):=(U+Z)(I+Z^\top Z)^{-1/2},\\
 R_U^{\mathrm{qr}}(Z):=(U+Z)\operatorname{chol}(I+Z^\top Z)^{-\top},\\
 R_U^{\mathrm{pol}}(0)=R_U^{\mathrm{qr}}(0)=U,\quad
 \mathrm DR_U^{\mathrm{pol}}(0)[Z]=\mathrm DR_U^{\mathrm{qr}}(0)[Z]=Z,\\
 \|R_U^{\mathrm{pol}}(Z)-U-Z\|_F\le
 \tfrac12\sqrt{1+\|Z\|_{\op}^2}\,\|Z\|_F^2,\\
 \|R_U^{\mathrm{qr}}(Z)-U-Z\|_F\le
 (1+\|Z\|_{\op}^2)\|Z\|_F^2.
 \label{eq:rigid-qr-error}
\end{gather}
\end{proposition}
\begin{proof}
\begin{align}
 X^\top X&=I+U^\top Z+Z^\top U+Z^\top Z=I+Z^\top Z,\\
 0&\preceq I-(I+E)^{-1/2}\preceq\tfrac12 E,\quad E:=Z^\top Z,\\
 \|R_U^{\mathrm{pol}}(Z)-X\|_F
 &\le\|X\|_{\op}\|(I+E)^{-1/2}-I\|_F
 \le\tfrac12\sqrt{1+\|Z\|_{\op}^2}\|E\|_F,\\
 L(s)L(s)^\top&=I+sE,\quad L(0)=I,\\
 L'(s)&=L(s)\mathcal T\bigl(L(s)^{-1}E L(s)^{-\top}\bigr),\\
 \mathcal T(A)&:=\operatorname{tril}(A,-1)+\tfrac12\diag(A),\\
 \|\mathcal T(A)\|_F&\le\|A\|_F,\quad
 \|L(s)^{-1}\|_{\op}\le1,\\
 \|L(1)-I\|_F&\le\int_0^1\|L'(s)\|_F\,\mathrm ds
 \le\sqrt{1+\|E\|_{\op}}\|E\|_F,\\
 \|XL(1)^{-\top}-X\|_F
 &\le\|X\|_{\op}\|L(1)^{-1}\|_{\op}\|L(1)-I\|_F\\
 &\le(1+\|E\|_{\op})\|E\|_F
 \le(1+\|Z\|_{\op}^2)\|Z\|_F^2.
\end{align}
\end{proof}
The retraction error is quadratic in the tangent step, so a small horizontal
update preserves its intended first-order motion. Comparing two different runs
additionally requires control of the QR map between their inputs. The next
lemma makes the necessary singular-value condition on that connecting segment
explicit.

\begin{lemma}[Differential of the positive-diagonal QR map]\label{lem:rigid-qr-lip}
For $X=QR$, $R_{ii}>0$, $\sigma_{\min}(X)\ge a>0$,
\begin{align}
 A&:=Q^\top(\mathrm dX)R^{-1},\\
 \mathrm dQ&=(I-QQ^\top)(\mathrm dX)R^{-1}
  +Q\bigl(\operatorname{tril}(A,-1)-\operatorname{tril}(A,-1)^\top\bigr),\\
 \|\mathrm dQ\|_F^2&\le2\|(\mathrm dX)R^{-1}\|_F^2
 \le\frac2{a^2}\|\mathrm dX\|_F^2.
\end{align}
If $\inf_{s\in[0,1]}\sigma_{\min}(X+s(Y-X))\ge a$,
\begin{equation}
 \|\operatorname{qf}(X)-\operatorname{qf}(Y)\|_F
 \le\frac{\sqrt2}{a}\|X-Y\|_F.
\end{equation}
\end{lemma}
\begin{proof}
\begin{align}
 \mathrm dX&=(\mathrm dQ)R+Q\,\mathrm dR,\\
 Q^\top(\mathrm dX)R^{-1}&=Q^\top\mathrm dQ+(\mathrm dR)R^{-1},\\
 (Q^\top\mathrm dQ)^\top&=-Q^\top\mathrm dQ,\qquad
 \operatorname{tril}((\mathrm dR)R^{-1},-1)=0,\\
 \|\operatorname{tril}(A,-1)-\operatorname{tril}(A,-1)^\top\|_F^2
 &=2\|\operatorname{tril}(A,-1)\|_F^2\le2\|A\|_F^2,\\
 \operatorname{qf}(Y)-\operatorname{qf}(X)
 &=\int_0^1\mathrm D\operatorname{qf}(X+s(Y-X))[Y-X]\,\mathrm ds.
\end{align}
\end{proof}

\subsection{Variational structure, spectral assignment, and local curvature}
Each prefix of the flag is evaluated against a covariance with its own memory
scale. Consequently a direction assigned to an inner block contributes to
several prefix objectives, as encoded by $A_j$. The assignment theorem below
resolves these competing contributions under simultaneous diagonalization,
while respecting the prescribed capacity of every block.

\begin{definition}[Flag energy]
\begin{gather}
 C_k=C_k^\top\succeq0,\quad \omega_k>0,\quad
 A_j:=\sum_{k=j}^K\omega_k C_k,\quad A_{K+1}:=0,\\
 \Phi(U;C):=\sum_{k=1}^K\omega_k\tr(P_k(U)C_k)
 =\sum_{j=1}^K\tr(U_j^\top A_jU_j),
 \label{eq:tf-objective}\\
 H(U;C):=2\sum_k\omega_k C_kUF_k,\quad
 \operatorname{grad}\Phi(U;C)=\mathsf H_U(H(U;C)).
\end{gather}
\end{definition}
\begin{proposition}[First and second variations]\label{prop:rigid-variation}
For $Z\in\Hor_U$, $Q=[U,U_\perp]$, and horizontal $\Omega$ with
$Z=(Q\Omega)_{:R}$,
\begin{align}
 \mathrm D\Phi(U)[Z]&=2\sum_k\omega_k\tr(Z^\top C_kUF_k)
 =\langle\operatorname{grad}\Phi,Z\rangle_F,\\
 \left.\frac{\mathrm d^2}{\mathrm ds^2}\Phi((Qe^{s\Omega})_{:R};C)\right|_{s=0}
 &=\sum_k\omega_k\tr\bigl(Q^\top C_kQ[\Omega,[\Omega,\widehat F_k]]\bigr),\\
 \operatorname{grad}\Phi(U;C)=0
 &\Longrightarrow\quad
 \langle Z,\operatorname{Hess}\Phi(U)[Z]\rangle_F
 =\sum_k\omega_k\tr\bigl(Q^\top C_kQ[\Omega,[\Omega,\widehat F_k]]\bigr).
\end{align}
\end{proposition}
\begin{proof}
\begin{align}
 \mathrm D(UF_kU^\top)[Z]&=ZF_kU^\top+UF_kZ^\top,\\
 \frac{\mathrm d}{\mathrm ds}(e^{s\Omega}\widehat F_ke^{-s\Omega})
 &=e^{s\Omega}[\Omega,\widehat F_k]e^{-s\Omega},\\
 \frac{\mathrm d^2}{\mathrm ds^2}(e^{s\Omega}\widehat F_ke^{-s\Omega})
 &=e^{s\Omega}[\Omega,[\Omega,\widehat F_k]]e^{-s\Omega},\\
 (\Phi\circ\gamma)''(0)
 &=\operatorname{Hess}\Phi[\dot\gamma(0),\dot\gamma(0)]
   +\langle\operatorname{grad}\Phi,\nabla_{\dot\gamma}\dot\gamma\rangle_F.
\end{align}
\end{proof}
\begin{theorem}[Commuting case: integral assignment]\label{thm:tf-assignment}
Suppose
\begin{gather}
 C_k=Q\diag(\lambda_{k1},\ldots,\lambda_{kd})Q^\top,\quad
 b_{ji}:=\sum_{k=j}^K\omega_k\lambda_{ki},\quad b_{K+1,i}:=0,\\
 \mathcal X:=\{x\ge0:\sum_{j=1}^{K+1}x_{ji}=1,
                    \ \sum_{i=1}^d x_{ji}=q_j\}.
\end{gather}
Then
\begin{align}
 \max_{[U]\in\mathfrak F}\Phi(U;C)
 &=\max_{x\in\mathcal X}\sum_{j,i}b_{ji}x_{ji}\\
 &=\max_{\substack{\{S_j\}\text{ partition }[d]\\ |S_j|=q_j}}
       \sum_{j=1}^{K+1}\sum_{i\in S_j}b_{ji}\\
 &=\min_{\alpha\in\RR^d,\,\beta\in\RR^{K+1}}
       \left\{\sum_i\alpha_i+\sum_jq_j\beta_j:
                    \alpha_i+\beta_j\ge b_{ji}\right\}.
\end{align}
\end{theorem}
\begin{proof}
\begin{align}
 x_{ji}(U)&:=e_i^\top Q^\top\Pi_j(U)Qe_i\ge0,\\
 \sum_jx_{ji}(U)&=1,\qquad \sum_ix_{ji}(U)=\tr\Pi_j=q_j,\\
 \Phi(U;C)&=\sum_{j,i}b_{ji}x_{ji}(U)
 \le\max_{x\in\mathcal X}\langle b,x\rangle.
\end{align}
Changing the signs of the block-sum rows turns the constraint matrix into a
directed bipartite incidence matrix. It is totally unimodular, so the
Hoffman--Kruskal integrality theorem~\citep{hoffman1956integral}, with integer
margins $(q_j,1)$, gives integral vertices. Column sums of one then give
$x_{ji}\in\{0,1\}$. For $S_j:=\{i:x_{ji}=1\}$ choose
$U_j=Q_{:,S_j}$. Then $x(U)=x$. Finally,
\begin{align}
 \mathcal L(x,\alpha,\beta)
 &=\sum_i\alpha_i+\sum_jq_j\beta_j+
       \sum_{j,i}(b_{ji}-\alpha_i-\beta_j)x_{ji},\\
 \sup_{x\ge0}\mathcal L<\infty
 &\Longleftrightarrow\alpha_i+\beta_j\ge b_{ji},\\
 \varnothing\ne\mathcal X\text{ compact}
 &\Longrightarrow\max\langle b,x\rangle=\min\mathcal L_{\mathrm{dual}}.
\end{align}
\end{proof}
\begin{corollary}[Common order and strict assignment margin]\label{cor:tf-common-order}
If $\lambda_{k1}\ge\cdots\ge\lambda_{kd}$ for every $k$, then
$U^\star=Q_{:,1:R}$ maximizes $\Phi$. More generally, let an optimal partition
$S^\star$ admit a dual optimum satisfying
\begin{gather}
 s_{ji}:=\alpha_i+\beta_j-b_{ji},\quad
 s_{ji}=0\ (i\in S_j^\star),\quad s_{ji}\ge\delta>0\ (i\notin S_j^\star),\\
 W(U):=\sum_j\sum_{i\notin S_j^\star}x_{ji}(U).
\end{gather}
Then
\begin{align}
 \Phi(U^\star)-\Phi(U)&=\sum_{j,i}s_{ji}x_{ji}(U)\ge\delta W(U),\\
 d_P(U,U^\star)^2&\le K W(U),\\
 \Phi(U^\star)-\Phi(U)&\ge\frac{\delta}{K}d_P(U,U^\star)^2.
\end{align}
The factor $K$ in the distance comparison is attained by a rotation between
the first and residual blocks.
\end{corollary}
\begin{proof}
Ky Fan's maximum principle~\citep{fan1949weyl} gives the common-order optimum.
For the quantitative bound let $T_{ja}:=\tr(\Pi_j(U)\Pi_a(U^\star))$, including
the residual block. Its row and column sums both equal the block capacities.
\begin{align}
 \tr(P_kC_k)&\le\sum_{i=1}^{r_k}\lambda_{ki}
 =\tr(P_k(U^\star)C_k),\\
 \|\Pi_j(U)-\Pi_j(U^\star)\|_F^2
 &=2q_j-2\sum_{i\in S_j^\star}x_{ji}
 =2\sum_{i\notin S_j^\star}x_{ji},\\
 \|P_k(U)-P_k(U^\star)\|_F^2
 &=2\sum_{j\le k<a}T_{ja}
 =\sum_{j\le k<a}T_{ja}+\sum_{a\le k<j}T_{ja},\\
 d_P^2&=\sum_{j,a}|j-a|T_{ja}\le K\sum_{j\ne a}T_{ja}=KW(U).
\end{align}
For a rotation of angle $\theta$ between a vector of block $1$ and one of
block $K+1$, $W=2\sin^2\theta$ and $d_P^2=2K\sin^2\theta$. Thus the geometric
constant cannot be reduced. The dual-slack identity gives the objective bound.
\end{proof}
\begin{proposition}[Perturbation and matched-rank reconstruction]
\label{prop:tf-value}
For $\|\widehat C_k-C_k\|_{\op}\le\epsilon_k$, define
$D_C:=\sum_k\omega_kr_k\epsilon_k$. If $U^\star,\widehat U$ maximize their
respective objectives, then
\begin{gather}
 \sup_U|\Phi(U;\widehat C)-\Phi(U;C)|\le D_C,\qquad
 0\le\Phi(U^\star;C)-\Phi(\widehat U;C)\le2D_C.
\end{gather}
For every rank-$R$ projector $Q_R$ containing $\operatorname{range}P_k$,
\begin{equation}\label{prop:tf-budget}
 \|G-GQ_R\|_F^2\le\|G-GP_k\|_F^2.
\end{equation}
\end{proposition}
\begin{proof}
\begin{align}
 |\tr(P_k\Delta C_k)|&\le\|P_k\|_*\|\Delta C_k\|_{\op}=r_k\epsilon_k,\\
 \Phi(U^\star;C)-\Phi(\widehat U;C)
 &\le\Phi(U^\star;C)-\Phi(U^\star;\widehat C)
       +\Phi(\widehat U;\widehat C)-\Phi(\widehat U;C),\\
 Q_R-P_k&\succeq0,\quad
 \|G-GP_k\|_F^2-\|G-GQ_R\|_F^2
 =\tr((Q_R-P_k)G^\top G)\ge0.
\end{align}
\end{proof}
\begin{proposition}[A strict nonglobal local maximum]\label{prop:tf-cycle}
Let
\begin{gather}
 d=3,\quad K=2,\quad(r_1,r_2)=(1,2),\quad\omega_1=\omega_2=1,\\
 C_1=\diag(3,\tfrac52,\tfrac12),\quad
 C_2=\diag(0,2,\tfrac32),\quad
 (b_{ji})=\begin{pmatrix}3&9/2&2\\0&2&3/2\\0&0&0\end{pmatrix},\\
 U_0=(e_1,e_2),\qquad U_\star=(e_2,e_3).
\end{gather}
Then $[C_1,C_2]=0$, $\Phi(U_0)=5<6=\Phi(U_\star)$, and
\begin{equation}
 \Phi((e^{\Omega})_{:2})=5-\tfrac12\Omega_{12}^2
                              -\Omega_{13}^2-\tfrac12\Omega_{23}^2
                              +O(\|\Omega\|_F^3).
\end{equation}
\end{proposition}
\begin{proof}
For eigen-indices $i\in S_a$, $l\in S_b$, $a<b$, put
\begin{gather}
 \Delta_{ai,bl}:=b_{ai}+b_{bl}-b_{al}-b_{bi},\\
 \left.\frac{\mathrm d^2}{\mathrm ds^2}\Phi((Qe^{s\Omega})_{:R})\right|_{0}
 =-2\sum_{a<b}\sum_{i\in S_a,l\in S_b}\Delta_{ai,bl}\Omega_{il}^2.
\end{gather}
At $U_0$,
\begin{gather}
 \Delta_{11,22}=\tfrac12,\quad \Delta_{11,33}=1,\quad
 \Delta_{22,33}=\tfrac12,\quad
 \operatorname{grad}\Phi(U_0)=0,\quad\operatorname{Hess}\Phi(U_0)\prec0.
\end{gather}
The six assignment scores are $5,9/2,9/2,6,2,4$.
\end{proof}
The counterexample distinguishes global allocation from local improvement:
every infinitesimal pairwise rotation at $U_0$ decreases the objective, although
a different assignment has greater value. Commutativity therefore provides an
exact global characterization without ensuring that ascent finds it. The
oracle used subsequently is an interior stationary branch with uniform local
curvature; the analysis requires the iterates to remain in its regular chart.

\begin{assumption}[A regular oracle chart]\label{ass:tf-chart}
Let $\psi:\mathcal D\subset\RR^p\to\mathfrak F$ be a smooth chart on a convex
set, $f_C:=\Phi(\psi(\cdot);C)$, and
\begin{gather}
 \nabla f_C(x_C^\star)=0,\quad
 -L_\phi I\preceq\nabla^2 f_C(x)\preceq-\mu_\phi I\quad(x\in\mathcal D),\\
 0<\mu_\phi\le L_\phi,\quad
 \|\mathrm D\psi(x)\|_{2\to F}\le J_\psi,\quad
 \|\nabla^3 f_C(x)\|\le T_\phi.
\end{gather}
\end{assumption}
\begin{lemma}[Oracle sensitivity and differentiability]\label{lem:tf-oracle}
If $C,\widehat C$ satisfy Assumption~\ref{ass:tf-chart} in the same chart,
then, with $E_C:=2J_\psi\sum_k\omega_k\sqrt{r_k}\epsilon_k$,
\begin{align}
 \sup_x\|\nabla f_{\widehat C}(x)-\nabla f_C(x)\|&\le E_C,\\
 \|x_{\widehat C}^\star-x_C^\star\|&\le E_C/\mu_\phi,\\
 \mathrm D x_C^\star[\Delta C]
 &=-[\nabla^2 f_C(x_C^\star)]^{-1}
                 \nabla f_{\Delta C}(x_C^\star).
\end{align}
If $\|\nabla^2f_{\Delta C}\|\le E_H$, and $v=x_{C+\Delta C}^\star-x_C^\star$,
\begin{equation}
 \|v-\mathrm D x_C^\star[\Delta C]\|
 \le\frac{T_\phi}{2\mu_\phi^3}E_C^2+
                         \frac{E_H}{\mu_\phi^2}E_C.
\end{equation}
\end{lemma}
\begin{proof}
\begin{align}
 \|H(U;\Delta C)\|_F&\le2\sum_k\omega_k\epsilon_k\|UF_k\|_F
 =2\sum_k\omega_k\sqrt{r_k}\epsilon_k,\\
 \mu_\phi\|v\|^2
 &\le-\langle\nabla f_C(x_C^\star+v)-\nabla f_C(x_C^\star),v\rangle
 \le E_C\|v\|,\\
 0&=\nabla f_{C+\Delta C}(x_C^\star+v)
 =\nabla^2 f_C(x_C^\star)v+\nabla f_{\Delta C}(x_C^\star)+R,\\
 \|R\|&\le\tfrac12T_\phi\|v\|^2+E_H\|v\|,\\
 v&=-[\nabla^2 f_C(x_C^\star)]^{-1}
                  (\nabla f_{\Delta C}(x_C^\star)+R).
\end{align}
\end{proof}
The factor $\mu_\phi^{-1}$ measures how strongly estimation errors can move the
local flag optimum. A nearly flat objective can change its preferred flag
substantially under a small covariance perturbation. The next spectral bound
exhibits the corresponding gap dependence for a single projector; a general
multi-covariance flag uses the chart curvature above.

\begin{proposition}[Residual Davis--Kahan bound]\label{prop:rigid-gap}
Let $C=Q_1\Lambda_1Q_1^\top+Q_2\Lambda_2Q_2^\top$, $\dim Q_1=r$, and
$g:=\lambda_r(C)-\lambda_{r+1}(C)>0$. If $\widehat C=C+E$,
$\|E\|_{\op}\le\epsilon<g$, and $\widehat P$ is its leading rank-$r$ projector,
set $s_r:=\min(r,d-r)$ and
$g_{\rm sep}:=\lambda_{\min}(\widehat\Lambda_1)-\lambda_{\max}(\Lambda_2)$.
Then $g_{\rm sep}\ge g-\epsilon>0$ and
\begin{equation}
 \|\widehat P-Q_1Q_1^\top\|_F
 \le\frac{\sqrt2\|Q_2^\top E\widehat Q_1\|_F}{g_{\rm sep}}
 \le\frac{\sqrt{2s_r}\,\epsilon}{g-\epsilon}.
\end{equation}
The left side is also at most $\sqrt{2s_r}$. Keeping the off-diagonal residual
can be substantially sharper than replacing it by $\sqrt{s_r}\epsilon$.
\end{proposition}
\begin{proof}
This is the Sylvester-equation form of the Davis--Kahan
$\sin\Theta$ theorem~\citep{davis_kahan}; Weyl's inequality supplies the
separation from the unperturbed gap.
\begin{align}
 \widehat C\widehat Q_1&=\widehat Q_1\widehat\Lambda_1,\quad
 X:=Q_2^\top\widehat Q_1,\quad
 X\widehat\Lambda_1-\Lambda_2X=Q_2^\top E\widehat Q_1=:B,\\
 \lambda_{\min}(\widehat\Lambda_1)-\lambda_{\max}(\Lambda_2)
 &\ge g-\epsilon,\\
 X_{ij}&=B_{ij}/(\widehat\lambda_j-\lambda_{r+i}),\quad
 \|X\|_F\le\|B\|_F/g_{\rm sep},\quad
 \|B\|_F\le\sqrt{s_r}\|E\|_{\op},\\
 \|\widehat P-P\|_F^2&=2r-2\tr(\widehat P P)=2\|Q_2^\top\widehat Q_1\|_F^2.
\end{align}
\end{proof}
\begin{corollary}[Local curvature at a common-order flag]
If $\lambda_{k1}\ge\cdots\ge\lambda_{kd}$, then for $i\in I_a$, $l\in I_b$,
\begin{gather}
 \Delta_{ai,bl}=\sum_{k=a}^{b-1}\omega_k(\lambda_{ki}-\lambda_{kl}),
 \quad \lambda_{k\cdot}:=0\ (k>K),\\
 \delta_\star:=\min_{a<b,\,i\in I_a,\,l\in I_b}\Delta_{ai,bl}>0
 \ \Longrightarrow\quad
 -\operatorname{Hess}\Phi(U^\star)[Z,Z]\ge\delta_\star\|Z\|_F^2.
\end{gather}
For orthonormal chart coordinates at $U^\star$, and
$\|\nabla^2f(x)-\nabla^2f(0)\|\le T_\phi\|x\|$,
\begin{equation}
 \|x\|\le\delta_\star/(2T_\phi)
 \quad\Longrightarrow\quad-\nabla^2f(x)\succeq(\delta_\star/2)I.
\end{equation}
\end{corollary}

\subsection{Adapted covariance statistics, sketches, and tracking}
Gradients depend on previous updates, so their second moments form an adapted
process. We separate fluctuations around each conditional covariance from the
lag introduced by averaging a changing target. Normalized and raw EMAs are
distinguished throughout because their early-time scaling changes the relative
weights in the flag objective.

\begin{definition}[Filtration and normalized EMA]
For a filtration $(\mathcal F_t)$, let $G_t$ be $\mathcal F_t$-measurable and
\begin{gather}
 \EE_t[\cdot]:=\EE[\cdot\mid\mathcal F_{t-1}],\\
 X_t:=G_t^\top G_t,\quad C_t:=\EE[X_t\mid\mathcal F_{t-1}],\\
 Z_t:=X_t-C_t,\quad\EE[Z_t\mid\mathcal F_{t-1}]=0,\\
 \overline C_{k,t}:=(1-\rho_k)\sum_{i=1}^t\rho_k^{t-i}X_i,\quad
 \widehat C_{k,t}:=\overline C_{k,t}/(1-\rho_k^t),\quad0<\rho_k<1,\\
 \alpha_{k,t,i}:=\frac{(1-\rho_k)\rho_k^{t-i}}{1-\rho_k^t},\qquad
 C^\circ_{k,t}:=\sum_{i=1}^t\alpha_{k,t,i}C_i.
\end{gather}
\end{definition}
\begin{lemma}[Exact EMA identities]\label{lem:rigid-ema}
\begin{align}
 \sum_{i=1}^t\alpha_{k,t,i}&=1,\quad
 \max_i\alpha_{k,t,i}=\frac{1-\rho_k}{1-\rho_k^t},\\
 \sum_{i=1}^t\alpha_{k,t,i}^2
 &=\frac{1-\rho_k}{1+\rho_k}\frac{1+\rho_k^t}{1-\rho_k^t},\\
 A_{k,t}:=\sum_i(t-i)\alpha_{k,t,i}
 &=\frac{\rho_k}{1-\rho_k}-\frac{t\rho_k^t}{1-\rho_k^t},\\
 \|C_i-C_{i-1}\|_{\op}\le\nu
 &\Longrightarrow\|C^\circ_{k,t}-C_t\|_{\op}\le\nu A_{k,t}.
\end{align}
\end{lemma}
\begin{proof}
\begin{align}
 \sum_{j=0}^{t-1}\rho^j&=(1-\rho^t)/(1-\rho),\quad
 \sum_{j=0}^{t-1}\rho^{2j}=(1-\rho^{2t})/(1-\rho^2),\\
 \sum_{j=0}^{t-1}j\rho^j
 &=\rho\frac{\mathrm d}{\mathrm d\rho}\frac{1-\rho^t}{1-\rho},\\
 \|C^\circ_{k,t}-C_t\|_{\op}
 &\le\sum_i\alpha_{k,t,i}\sum_{j=i+1}^t\|C_j-C_{j-1}\|_{\op}
 \le\nu\sum_i(t-i)\alpha_{k,t,i}.
\end{align}
\end{proof}
\begin{theorem}[Finite-horizon adapted covariance bound]\label{thm:rigid-covariance}
Suppose, almost surely,
\begin{gather}
 \|Z_i\|_{\op}\le L_Z,\quad
 \EE[Z_i^2\mid\mathcal F_{i-1}]\preceq v_i I,
 \quad v_i\text{ deterministic},\\
 V_{k,t}:=\sum_i\alpha_{k,t,i}^2v_i,\quad
 R_{k,t}:=L_Z\max_i\alpha_{k,t,i},\\
 u:=\log(2dKT/\delta),\quad
 \mathfrak r(V,R,u):=\frac{Ru}{3}+\sqrt{2Vu+\frac{R^2u^2}{9}}.
\end{gather}
With probability at least $1-\delta$, simultaneously for $k\le K,t\le T$,
\begin{align}
 \|\widehat C_{k,t}-C_t\|_{\op}
 &\le\mathfrak r(V_{k,t},R_{k,t},u)+\nu A_{k,t},\\
 \|\overline C_{k,t}-C_t\|_{\op}
 &\le(1-\rho_k^t)
  \left(\mathfrak r(V_{k,t},R_{k,t},u)+\nu A_{k,t}\right)
       +\rho_k^t\|C_t\|_{\op}.
\end{align}
\end{theorem}
\begin{proof}
For fixed $(k,t)$, $D_i:=\alpha_{k,t,i}Z_i$ satisfies
\begin{gather}
 \EE[D_i\mid\mathcal F_{i-1}]=0,\quad
 \lambda_{\max}(\pm D_i)\le R_{k,t},\quad
 \left\|\sum_i\EE[D_i^2\mid\mathcal F_{i-1}]\right\|_{\op}\le V_{k,t}.
\end{gather}
The matrix Freedman inequality \citep{tropp2011freedman} gives
\begin{align}
 \PP\left(\left\|\sum_iD_i\right\|_{\op}\ge s\right)
 &\le2d\exp\left(-\frac{s^2}{2(V_{k,t}+R_{k,t}s/3)}\right),\\
 s=\mathfrak r(V_{k,t},R_{k,t},u)
 &\Longrightarrow s^2=2u(V_{k,t}+R_{k,t}s/3),\\
 \PP\left(\exists k,t:\|\widehat C_{k,t}-C^\circ_{k,t}\|_{\op}>s\right)
 &\le2dKT e^{-u}=\delta,\\
 \overline C_{k,t}-C_t
 &=(1-\rho_k^t)(\widehat C_{k,t}-C_t)-\rho_k^t C_t.
\end{align}
Thus the radius exactly inverts the Bernstein exponent in matrix Freedman;
it avoids the further relaxation
$\mathfrak r(V,R,u)\le\sqrt{2Vu}+2Ru/3$. A zero conditional-variance budget
means the martingale sum vanishes almost surely and its radius can be set to zero.
\end{proof}
\begin{corollary}[Memory--drift balance]
For the limiting EMA weights, $v_i\le v$, $h=1-\rho$, and fixed $u$,
\begin{gather}
 \epsilon(h)\le\mathfrak r\left(\frac{vh}{2-h},L_Zh,u\right)
                         +\frac{\nu(1-h)}h,\\
 \mathfrak r\left(\frac{vh}{2-h},L_Zh,u\right)
 =\sqrt{vu}\sqrt h+\frac{L_Zu}{3}h+O(h^{3/2}),\quad h\downarrow0,\\
 a:=\sqrt{vu}>0,\quad\nu>0,\quad
 h_0:=(2\nu/a)^{2/3}\le1,\quad
 a\sqrt{h_0}+\nu/h_0=3\,2^{-2/3}a^{2/3}\nu^{1/3}.
\end{gather}
The choice $h_0$ minimizes the leading stochastic--drift expression
$a\sqrt h+\nu/h$; the first display retains the finite-$h$ constants.
\end{corollary}
Longer memory reduces the stochastic term by spreading mass over more
observations, but increases the age of the information used to track the
current covariance. The displayed balance quantifies this tradeoff. Different
timescales can address different drift regimes; their joint benefit still
depends on the flag objective and the resulting update.

\begin{definition}[Fixed-probe Nystr\"om reconstruction]
\begin{gather}
 \Omega_k\in\RR^{d\times s_k},\quad s_k=\min(d,r_k+p),\quad
 Y_k=C_k\Omega_k,\quad B_k=\sym(\Omega_k^\top Y_k)+\gamma I,\quad\gamma>0,\\
 \mathcal N_\gamma(Y_k):=Y_kB_k^{-1}Y_k^\top.
\end{gather}
\end{definition}
\begin{proposition}[PSD order, regularization, and probe blindness]\label{prop:rigid-nystrom}
For $C\succeq0$, $Y=C\Omega$, and $Z=C^{1/2}\Omega$,
\begin{align}
 \mathcal N_\gamma(C\Omega)
 &=C^{1/2}Z(Z^\top Z+\gamma I)^{-1}Z^\top C^{1/2},\\
 0\preceq\mathcal N_\gamma(C\Omega)&\preceq C,\\
 \mathcal N_0(C\Omega)-\mathcal N_\gamma(C\Omega)
 &=C^{1/2}L\diag\left(\frac{\gamma}{\sigma_i^2+\gamma}\right)L^\top C^{1/2},
 \quad Z=L\diag(\sigma_i)R^\top\quad(\sigma_i>0).
\end{align}
If $s<d$, $v\in\ker\Omega^\top$, and $\|v\|=1$, then
\begin{equation}
 C=a vv^\top\quad\Longrightarrow\quad Y=0,\quad
 \mathcal N_\gamma(Y)=0,\quad\|C-\mathcal N_\gamma(Y)\|_{\op}=a.
\end{equation}
Here $\mathcal N_0$ uses the Moore--Penrose inverse.
\end{proposition}
\begin{proof}
\begin{align}
 Z(Z^\top Z+\gamma I)^{-1}Z^\top
 &=L\diag\left(\frac{\sigma_i^2}{\sigma_i^2+\gamma}\right)L^\top,\\
 0&\le\frac{\sigma_i^2}{\sigma_i^2+\gamma}\le1,\quad
 Z(Z^\top Z)^\dagger Z^\top=LL^\top,\\
 C\Omega&=av(v^\top\Omega)=0.
\end{align}
\end{proof}
\begin{lemma}[Sketch sensitivity]\label{lem:rigid-sketch-lip}
Suppose $B=\sym(\Omega^\top Y)+\gamma I\succeq\gamma I$ and
$B'=\sym(\Omega^\top Y')+\gamma I\succeq\gamma I$, with
$\max(\|Y\|_{\op},\|Y'\|_{\op})\le Y_\star$. Then
\begin{align}
 \|\mathcal N_\gamma(Y)-\mathcal N_\gamma(Y')\|_F
 &\le L_N\|Y-Y'\|_F,\quad
 L_N:=2Y_\star/\gamma+Y_\star^2\|\Omega\|_{\op}/\gamma^2,\\
 \mathrm D\mathcal N_\gamma(Y)[E]
 &=EB^{-1}Y^\top+YB^{-1}E^\top
                   -YB^{-1}\sym(\Omega^\top E)B^{-1}Y^\top.
\end{align}
\end{lemma}
\begin{proof}
\begin{align}
 B^{-1}-B'^{-1}&=B^{-1}(B'-B)B'^{-1},\\
 YB^{-1}Y^\top-Y'B'^{-1}Y'^\top
 &=(Y-Y')B^{-1}Y^\top+Y'B^{-1}(Y-Y')^\top
                       +Y'(B^{-1}-B'^{-1})Y'^\top,\\
 \|B^{-1}-B'^{-1}\|_F
 &\le\gamma^{-2}\|\Omega\|_{\op}\|Y-Y'\|_F,\qquad
 \mathrm D(B^{-1})[E]=-B^{-1}\sym(\Omega^\top E)B^{-1}.
\end{align}
\end{proof}
The probe-blindness example explains why sketch error must remain a separate
quantity: a fixed probe can miss an entire energetic direction. We now combine
this approximation error with statistical error and target motion. The
tracking recursion measures distance to a local oracle at refresh times, rather
than assuming exact covariance information at every parameter update.

\begin{definition}[Refresh-time chart recursion]
At refresh index $s$, let
\begin{gather}
 f_s(x):=\Phi(\psi(x);C_s),\quad \nabla f_s(x_s^\star)=0,\\
 x_{s+1}=x_s+\xi\nabla f_s(x_s)+\xi e_s+\xi\zeta_s+r_s,\\
 \Delta_s:=x_{s+1}^\star-x_s^\star,\quad
 E_s:=\EE\|x_s-x_s^\star\|^2,\\
 q:=\max\{|1-\xi\mu_\phi|,|1-\xi L_\phi|\}<1,\quad0<\xi<2/L_\phi.
\end{gather}
\end{definition}
\begin{theorem}[Cocoercive tracking with exact forcing gains]\label{thm:tf-tracking}
Under Assumption~\ref{ass:tf-chart}, assume the relevant iterates and connecting
segments remain in its convex chart. Define
\begin{gather}
 v_s:=\xi e_s+\xi\zeta_s+r_s-\Delta_s,\quad
 b_s:=\sqrt{\EE\|v_s\|^2}.
\end{gather}
Then the root-mean-square error obeys
\begin{align}
 \sqrt{E_{s+1}}&\le q\sqrt{E_s}+b_s,\\
 \sqrt{E_s}&\le q^s\sqrt{E_0}+\sum_{j<s}q^{s-1-j}b_j.
 \label{eq:tracking-rms-sharp}
\end{align}
The contraction is minimized at
$\xi_\star=2/(L_\phi+\mu_\phi)$, giving
$q_\star=(L_\phi-\mu_\phi)/(L_\phi+\mu_\phi)$.
For comparison, a squared-error form remains valid: for $\alpha>0$ with
$\kappa:=(1+\alpha)q^2<1$,
\begin{align}
 E_{s+1}&\le\kappa E_s+B_s,\\
 B_s&:=4(1+\alpha^{-1})\EE\left[
       \xi^2\|e_s\|^2+\xi^2\|\zeta_s\|^2+\|r_s\|^2+\|\Delta_s\|^2\right],\\
 E_s&\le\kappa^s E_0+\sum_{j=0}^{s-1}\kappa^{s-1-j}B_j,
 \label{eq:tf-tracking-expanded}\\
 \sum_{s=0}^{S-1}E_s&\le\frac{E_0+\sum_{s=0}^{S-2}B_s}{1-\kappa}.
\end{align}
If $\EE[\zeta_s\mid\mathcal H_s]=0$, where $\mathcal H_s$ contains
$(f_s,x_s,x_s^\star,e_s,r_s,\Delta_s)$, set
$d_s^2:=\EE\|\xi e_s+r_s-\Delta_s\|^2$ and
$\sigma_s^2:=\EE\|\zeta_s\|^2$. Then
\begin{equation}
 E_{s+1}\le(q\sqrt{E_s}+d_s)^2+\xi^2\sigma_s^2.
\end{equation}
For uniform envelopes $d_s\le\overline d$, $\sigma_s\le\overline\sigma$, define
\begin{equation}
 x_\infty:=\frac{q\overline d+
      \sqrt{\overline d^2+(1-q^2)\xi^2\overline\sigma^2}}{1-q^2}.
\end{equation}
Then $\sqrt{E_s}\le x_\infty+q^s(\sqrt{E_0}-x_\infty)_+$.
In particular the pure-drift RMS gain is $1/(1-q)$, whereas the centered-noise
variance gain is $\xi^2/(1-q^2)$; scalar quadratic recursions attain both.
\end{theorem}
\begin{proof}
The smooth strongly convex interpolation inequality associated with
Baillon--Haddad cocoercivity~\citep{bauschke2010baillon} gives the optimal-step
contraction. Here its local Hessian proof also gives every admissible step:
with $v=x-y$ and $g=\nabla(-f_s)(x)-\nabla(-f_s)(y)$,
\begin{align}
 \langle v,g\rangle&\ge\frac{\mu_\phi L_\phi}{\mu_\phi+L_\phi}\|v\|^2
                 +\frac{\|g\|^2}{\mu_\phi+L_\phi},\\
 T_s(x)-T_s(y)
 &=\left[I+\xi\int_0^1\nabla^2f_s(y+u(x-y))\,\mathrm du\right](x-y),\\
 \|T_s(x)-T_s(y)\|&\le q\|x-y\|,\quad T_s(x_s^\star)=x_s^\star,\\
 x_{s+1}-x_{s+1}^\star
 &=T_s(x_s)-T_s(x_s^\star)+\xi e_s+\xi\zeta_s+r_s-\Delta_s,\\
 \|a+b\|^2&\le(1+\alpha)\|a\|^2+(1+\alpha^{-1})\|b\|^2,\\
 \|b_1+b_2+b_3+b_4\|^2&\le4\sum_{i=1}^4\|b_i\|^2,\\
 \sum_{s=0}^{S-1}\sum_{j<s}\kappa^{s-1-j}B_j
 &\le(1-\kappa)^{-1}\sum_{j=0}^{S-2}B_j.
\end{align}
The averaged Hessian in the first map identity has spectrum in
$[-L_\phi,-\mu_\phi]$, so the endpoint formula for $q$ is exact for the
quadratic class. Minkowski's inequality in $L^2$ gives
\eqref{eq:tracking-rms-sharp} directly. Conditional centering removes the
cross term with $\zeta_s$. The scalar map
$x\mapsto\sqrt{(qx+\overline d)^2+\xi^2\overline\sigma^2}$ has Lipschitz constant at most
$q$ and fixed point $x_\infty$, proving the last bound. Constant scalar forcing
and a scalar AR(1) process respectively attain the two stated gains.
Sharpness here concerns the comparison recursion; applicability to a flag
trajectory still requires the stated chart invariance.
\end{proof}
\begin{corollary}[End-to-end flag forcing]\label{cor:rigid-flag-forcing}
Suppose $\|\widetilde C_{k,s}-\widehat C_{k,s}\|_{\op}\le a_{k,s}$ and
$\|\widehat C_{k,s}-C_{k,s}\|_{\op}\le\epsilon_{k,s}$. Then
\begin{align}
 \|e_s\|&\le2J_\psi\sum_k\omega_k\sqrt{r_k}(a_{k,s}+\epsilon_{k,s}),\\
 \|\Delta_s\|&\le\frac{2J_\psi}{\mu_\phi}
               \sum_k\omega_k\sqrt{r_k}\|C_{k,s+1}-C_{k,s}\|_{\op}.
\end{align}
If $\|C_k(w)-C_k(w')\|_{\op}\le L_{C,k}\|w-w'\|$ and
$w_{t+1}=w_t-\eta_t d_t$, with $\tau$ steps per refresh, then
\begin{align}
 \|\Delta_s\|^2
 &\le L_\star^2\tau\sum_{t=s\tau}^{(s+1)\tau-1}\eta_t^2\|d_t\|^2,
 \quad L_\star:=\frac{2J_\psi}{\mu_\phi}\sum_k\omega_k\sqrt{r_k}L_{C,k}.
\end{align}
\end{corollary}
\begin{proposition}[Bounded-gradient constants and raw-EMA weights]\label{prop:rigid-raw-ema}
If $\|G_t\|_F\le G_\star$ almost surely, then
\begin{gather}
 0\preceq X_t,C_t\preceq G_\star^2I,\quad
 \|Z_t\|_{\op}\le G_\star^2,\quad
 \EE_tZ_t^2=\EE_tX_t^2-C_t^2\preceq G_\star^4I,\\
 L_Z=G_\star^2,\qquad v_t=G_\star^4
 \quad\text{are admissible in Theorem~\ref{thm:rigid-covariance}},\\
 \Phi(U;\overline C_{1:K,t})
 =\sum_k\omega_k(1-\rho_k^t)\tr(P_k(U)\widehat C_{k,t}).
\end{gather}
Thus replacing raw EMAs by normalized EMAs changes the relative objective weights
unless all factors $1-\rho_k^t$ coincide.
\end{proposition}
\begin{proof}
\begin{gather}
 -G_\star^2I\preceq X_t-C_t\preceq G_\star^2I,\quad
 X_t^2\preceq G_\star^2X_t,\\
 \EE_t(X_t-C_t)^2=\EE_tX_t^2-C_t^2
 \preceq G_\star^2 C_t\preceq G_\star^4I.
\end{gather}
\end{proof}
The preceding chart recursion uses a Euclidean gradient for clarity. Actual
manifold ascent uses the inverse pullback metric, and QR introduces a retraction
remainder. The next theorem controls both effects so that the same tracking
argument applies under explicit local metric and step-size conditions.

\begin{theorem}[Actual Riemannian ascent in a non-Euclidean chart]
\label{thm:rigid-chart-metric}
Under Assumption~\ref{ass:tf-chart}, let $\mathcal D$ lie in the radius-$r$ ball
about $x^\star$, and define the pullback metric $\mathsf G(x)$ and natural gradient
\begin{gather}
 \mathsf G(x)_{ab}:=g_{\psi(x)}(\mathrm D\psi(x)e_a,\mathrm D\psi(x)e_b),\quad
 A(x):=\mathsf G(x)^{-1}\nabla f(x),\\
 \|\mathsf G^{-1}(x)-I\|_{\op}\le\delta_M,\quad
 \|\mathrm D(\mathsf G^{-1})(x)\|\le L_M,\\
 \mu_A:=\mu_\phi-L_\phi(\delta_M+L_Mr)>0,\quad
 L_A:=L_\phi(1+\delta_M+L_Mr),\\
 0<\xi<2\mu_A/L_A^2,\quad q_A:=\sqrt{1-2\mu_A\xi+L_A^2\xi^2}<1.
\end{gather}
Then
\begin{align}
 \langle x-y,A(x)-A(y)\rangle&\le-\mu_A\|x-y\|^2,\\
 \|A(x)-A(y)\|&\le L_A\|x-y\|,\\
 \|x+\xi A(x)-y-\xi A(y)\|&\le q_A\|x-y\|.
\end{align}
Let $\mathcal R_x(v)$ be the chart representation of the quotient QR retraction,
with $\mathcal R_x(0)=x$, $\mathrm D\mathcal R_x(0)=I$, and
$\|\mathrm D^2\mathcal R_x(v)\|\le C_\psi$ on the relevant segment. For an
estimated covariance $\widehat C$, the actual ascent is
\begin{gather}
 x^+=\mathcal R_x\left(\xi\mathsf G(x)^{-1}\nabla f_{\widehat C}(x)\right)
 =x+\xi A(x)+\xi e(x)+r(x),\\
 \|e(x)\|\le(1+\delta_M)E_C,\qquad
 \|r(x)\|\le\tfrac12 C_\psi\xi^2\|A(x)+e(x)\|^2.
\end{gather}
Theorem~\ref{thm:tf-tracking} therefore applies with $q$ replaced by $q_A$,
$\kappa=(1+\alpha)q_A^2<1$, and the displayed error and remainder bounds.
\end{theorem}
\begin{proof}
For $M(x):=\mathsf G(x)^{-1}$,
\begin{align}
 \mathrm DA(x)[v]&=\nabla^2f(x)v+(M(x)-I)\nabla^2f(x)v
                           +(\mathrm DM(x)[v])\nabla f(x),\\
 \|\nabla f(x)\|&=\|\nabla f(x)-\nabla f(x^\star)\|\le L_\phi r,\\
 v^\top\mathrm DA(x)v
 &\le-\mu_\phi\|v\|^2+\delta_M L_\phi\|v\|^2+L_ML_\phi r\|v\|^2
 =-\mu_A\|v\|^2,\\
 \|\mathrm DA(x)\|_{\op}&\le L_A,\\
 \|x-y+\xi(A(x)-A(y))\|^2
 &\le(1-2\mu_A\xi+L_A^2\xi^2)\|x-y\|^2,\\
 e(x)&:=M(x)(\nabla f_{\widehat C}(x)-\nabla f_C(x)),\\
 r(x)&=\int_0^1(1-s)\mathrm D^2\mathcal R_x(sv)[v,v]\,\mathrm ds,
 \quad v=\xi(A(x)+e(x)).
\end{align}
\end{proof}

\subsection{The implemented adaptive map and its coordinate errors}
This subsection fixes the order of operations in the implemented optimizer.
The parameter step uses the current flag; a scheduled refresh then aligns the
new basis, transports first moments, and resets second moments. Keeping their
ages separate is essential because the two statistics retain different
amounts of history after a refresh.

\begin{definition}[One exact matrix-state transition]\label{def:rigid-map}
Let $W_t,G_t\in\RR^{m\times d}$, $U_t\in\St(d,R)$, and
$\chi_t:=\mathbf1_{t\equiv0\pmod\tau}$. Operations $\odot$, division, and square
roots below are entrywise. Assume $0\le\beta_{jk}<1$, $\epsilon,\gamma,\xi>0$,
$\tau\in\mathbb N$. Before the refresh,
\begin{align}
 Y_{k,t}^-&=\rho_kY_{k,t-1}+(1-\rho_k)G_t^\top G_t\Omega_k,\\
 Z_{k,t}&=G_tU_{k,t},\\
 M_{k,t}^-&=\beta_{1k}M_{k,t-1}^++(1-\beta_{1k})Z_{k,t},\\
 V_{k,t}^-&=\beta_{2k}V_{k,t-1}^++(1-\beta_{2k})Z_{k,t}^{\odot2},\\
 c_{M,k,t}^-&=c_{M,k,t-1}^++1,\quad c_{V,k,t}^-=c_{V,k,t-1}^++1,\\
 \widehat M_{k,t}&=M_{k,t}^-/(1-\beta_{1k}^{c_{M,k,t}^-}),\quad
 \widehat V_{k,t}=V_{k,t}^-/(1-\beta_{2k}^{c_{V,k,t}^-}),\\
 H_{k,t}&=\widehat M_{k,t}/(\sqrt{\widehat V_{k,t}}+\epsilon),\\
 D_t&=\sum_k a_kH_{k,t}U_{k,t}^\top+
   a_\perp\mathcal R_\epsilon(G_t(I-U_tU_t^\top)),
 \label{eq:rigid-direction}\\
 \mathcal R_\epsilon(X)&:=X/(\|X\|_F/\sqrt{md}+\epsilon),\\
 W_{t+1}&=(1-\eta_t\lambda)W_t-\eta_tD_t.
\end{align}
If $\chi_t=1$,
\begin{align}
 \widetilde C_{k,t}&:=\mathcal N_\gamma(Y_{k,t}^-),\quad
 \overline U_t:=\operatorname{qf}(U_t+\xi\operatorname{grad}\Phi(U_t;\widetilde C_t)),\\
 O_{k,t}&:=\operatorname{polar}(\overline U_{k,t}^\top U_{k,t}),\quad
 U_{k,t+1}:=\overline U_{k,t}O_{k,t},\\
 M_{k,t}^+&:=M_{k,t}^-O_{k,t}^\top,\quad
 V_{k,t}^+:=0,\quad c_{V,k,t}^+:=0,\quad c_{M,k,t}^+:=c_{M,k,t}^-.
\end{align}
If $\chi_t=0$, set $(U_{t+1},M_t^+,V_t^+,c_t^+)=(U_t,M_t^-,V_t^-,c_t^-)$. In both
cases $Y_{k,t}=Y_{k,t}^-$. Thus $D_t$ uses $U_t$, not $U_{t+1}$.
\end{definition}
\begin{lemma}[Procrustes optimization and polar sensitivity]\label{lem:rigid-polar}
For orthonormal $A,B\in\RR^{d\times q}$ and nonsingular $A^\top B=L\Sigma R^\top$,
\begin{gather}
 \arg\min_{O\in\mathrm O(q)}\|AO-B\|_F^2=LR^\top,\qquad
 \min_O\|AO-B\|_F^2=2q-2\|A^\top B\|_*,\\
 \|\operatorname{polar}(X)-\operatorname{polar}(Y)\|_F
 \le\frac{2\|X-Y\|_F}{\sigma_{\min}(X)+\sigma_{\min}(Y)}.
\end{gather}
\end{lemma}
\begin{proof}
\begin{align}
 \|AO-B\|_F^2&=2q-2\tr(O^\top A^\top B),\\
 \tr(O^\top L\Sigma R^\top)&=\tr(\Sigma R^\top O^\top L)
 \le\tr\Sigma,\\
 X&=PH,\quad Y=P'H',\quad Q=P^\top P',\quad E=X-Y,\\
 P^\top E-E^\top P'&=H(I-Q)+(I-Q)H',\\
 \langle T,HT+TH'\rangle_F
 &\ge(\lambda_{\min}(H)+\lambda_{\min}(H'))\|T\|_F^2,\\
 \|I-Q\|_F&\le\frac{\|P^\top E-E^\top P'\|_F}
                         {\lambda_{\min}(H)+\lambda_{\min}(H')}
 \le\frac{2\|E\|_F}{\sigma_{\min}(X)+\sigma_{\min}(Y)}.
\end{align}
\end{proof}
\begin{proposition}[Exact least-squares transport and implemented discrepancy]
\label{prop:rigid-transport}
For old and new block bases $U,V\in\St(d,q)$, $A:=MU^\top$, the unique minimizer is
\begin{align}
 M_\mathrm{LS}&:=\arg\min_N\|NV^\top-A\|_F^2=MU^\top V,\\
 \|NV^\top-A\|_F^2
 &=\|N-MU^\top V\|_F^2+\|MU^\top(I-VV^\top)\|_F^2.
\end{align}
For $V=\overline U O$ and the implemented $M^+=MO^\top$,
\begin{align}
 \|M^+V^\top-MU^\top\|_F^2
 &=\|M(O^\top-U^\top V)\|_F^2+
                              \|MU^\top(I-VV^\top)\|_F^2,\\
 \|MU^\top(I-VV^\top)\|_F
 &\le\|M\|_{\op}\|UU^\top-VV^\top\|_F/\sqrt2.
\end{align}
In particular,
\begin{gather}
 \overline U=UR,\quad O=R^\top,\quad V=U,\quad M^+=MR,\\
 d_P(U,V)=0,\qquad \|M^+V^\top-MU^\top\|_F=\|M(R-I)\|_F.
\end{gather}
\end{proposition}
\begin{proof}
\begin{align}
 NV^\top-A&=(N-AV)V^\top-A(I-VV^\top),\\
 \langle(N-AV)V^\top,A(I-VV^\top)\rangle_F&=0,\\
 \|U^\top(I-VV^\top)\|_F^2
 &=q-\tr(UU^\top VV^\top)=\tfrac12\|UU^\top-VV^\top\|_F^2,\\
 \operatorname{polar}((UR)^\top U)&=R^\top.
\end{align}
\end{proof}
The transport identity separates two losses of information. Projection onto a
changed subspace discards the component outside it, while a coordinate mismatch
can perturb the represented moment even inside the new subspace. Least-squares
transport minimizes their total squared error. The implemented rotation has
an additional, explicitly retained discrepancy, including in the pure-basis
example above.

\begin{lemma}[First-moment unrolling and boundedness]\label{lem:rigid-moment}
If $M_0=0$, $M_t^-=\beta M_{t-1}^++(1-\beta)Z_t$, $M_t^+=M_t^-T_t$, and
$\|T_t\|_{\op}\le1$, then
\begin{align}
 M_t^-&=(1-\beta)\sum_{i=1}^t\beta^{t-i}Z_iT_iT_{i+1}\cdots T_{t-1},\\
 \|Z_i\|_F\le G_\star
 &\Longrightarrow\|M_t^-\|_F\le(1-\beta^t)G_\star,\quad
 \|\widehat M_t\|_F\le G_\star.
\end{align}
For $T_t=O_t^\top$ or $T_t=U_t^\top U_{t+1}$, $\|T_t\|_{\op}\le1$. If the second
moment age is $a_t\ge1$ and $\|Z_i\|_F\le G_\star$, then
\begin{gather}
 \|\sqrt{\widehat V_t}\|_F^2\le G_\star^2,\quad
 \|H_t\|_F\le G_\star/\epsilon,\quad
 \|D_t\|_F\le\epsilon^{-1}G_\star\sum_k|a_k|+|a_\perp|\sqrt{md}.
\end{gather}
\end{lemma}
\begin{proof}
\begin{align}
 \|M_t^-\|_F&\le(1-\beta)G_\star\sum_{i=1}^t\beta^{t-i}
 =(1-\beta^t)G_\star,\\
 \tr\widehat V_t&=\frac{1-\beta_2}{1-\beta_2^{a_t}}
                 \sum_{j=0}^{a_t-1}\beta_2^j\|Z_{t-j}\|_F^2\le G_\star^2,\\
 \|H_t\|_F&\le\|\widehat M_t\|_F/\epsilon,\quad
 \|\mathcal R_\epsilon(X)\|_F
 =\frac{\|X\|_F}{\|X\|_F/\sqrt{md}+\epsilon}\le\sqrt{md}.
\end{align}
\end{proof}
\begin{lemma}[Square-root state and normalization]\label{lem:rigid-root}
For $S=\sqrt V$, define
$\mathcal S_\beta(S,Z):=\sqrt{\beta S^{\odot2}+(1-\beta)Z^{\odot2}}$. Then
\begin{align}
 \|\mathcal S_\beta(S,Z)-\mathcal S_\beta(S',Z')\|_F
 &\le\sqrt\beta\|S-S'\|_F+\sqrt{1-\beta}\|Z-Z'\|_F,\\
 \left\|\frac{M}{S+\epsilon}-\frac{M'}{S'+\epsilon}\right\|_F
 &\le\epsilon^{-1}\|M-M'\|_F+
          \epsilon^{-2}\|M'\|_{\max}\|S-S'\|_F.
\end{align}
For synchronized ages $c\ge1,a\ge1$ and $\|M'\|_{\max}\le G_\star$,
\begin{align}
 \|H-H'\|_F&\le A_M\|M-M'\|_F+A_S\|S-S'\|_F,\\
 A_M&:=\frac1{(1-\beta_1)\epsilon},\quad
 A_S:=\frac{G_\star}{(1-\beta_1)\epsilon^2\sqrt{1-\beta_2}}.
\end{align}
\end{lemma}
\begin{proof}
\begin{align}
 |\mathcal S_\beta(s,z)-\mathcal S_\beta(s',z')|
 &=\big|\|(\sqrt\beta s,\sqrt{1-\beta}z)\|_2
       -\|(\sqrt\beta s',\sqrt{1-\beta}z')\|_2\big|\\
 &\le\sqrt{\beta|s-s'|^2+(1-\beta)|z-z'|^2},\\
 \frac{m}{s+\epsilon}-\frac{m'}{s'+\epsilon}
 &=\frac{m-m'}{s+\epsilon}
   +\frac{m'(s'-s)}{(s+\epsilon)(s'+\epsilon)}.
\end{align}
\end{proof}
Using $\sqrt V$ as the state distance avoids differentiating the square-root
function at a zero variance entry. The resulting normalization bound remains
finite because of $\epsilon$, but its sensitivity can scale as
$\epsilon^{-2}$. This identifies a concrete connection between numerical
regularization and the later stability constants.

\begin{proposition}[Radial residual map]\label{prop:rigid-radial}
Writing $x=\operatorname{vec}X$, $N=md$, $r=\|x\|$, the eigenvalues of
$\mathrm D\mathcal R_\epsilon(x)$ are
\begin{gather}
 \lambda_{\mathrm{rad}}=\frac{\epsilon}{(r/\sqrt N+\epsilon)^2},\quad
 \lambda_{\mathrm{tan}}=\frac1{r/\sqrt N+\epsilon},\\
 0\preceq\mathrm D\mathcal R_\epsilon(x)\preceq\epsilon^{-1}I,\quad
 \|\mathcal R_\epsilon(x)-\mathcal R_\epsilon(y)\|\le\epsilon^{-1}\|x-y\|,\\
 \langle x,\mathcal R_\epsilon(x)\rangle
 =\frac{r^2}{r/\sqrt N+\epsilon}\ge0.
\end{gather}
At $x=0$, $\mathrm D\mathcal R_\epsilon(0)=\epsilon^{-1}I$.
\end{proposition}
\begin{proof}
\begin{equation}
 \mathrm D\mathcal R_\epsilon(x)
 =\frac{I}{r/\sqrt N+\epsilon}
    -\frac{xx^\top}{\sqrt N\,r(r/\sqrt N+\epsilon)^2}\quad(r>0).
\end{equation}
\end{proof}
\begin{proposition}[Coordinatewise Adam is not a quotient map]\label{prop:tf-gauge}
For $U=(e_1,e_2)\in\St(3,2)$, $G=(1,0,0)$, zero initial moments, and
$R=2^{-1/2}\left(\begin{smallmatrix}1&-1\\1&1\end{smallmatrix}\right)$,
\begin{gather}
 UU^\top=(UR)(UR)^\top,\\
 \frac{GU}{|GU|+\epsilon}U^\top=\frac{e_1^\top}{1+\epsilon},\qquad
 \frac{GUR}{|GUR|+\epsilon}R^\top U^\top
                 =\frac{e_1^\top}{2^{-1/2}+\epsilon}.
\end{gather}
For a signed permutation $Q$, the transformations
\begin{gather}
 U'=UQ,\quad Z'=ZQ,\quad M'=MQ,\quad V'=V|Q|,\\
 H'=HQ,\quad H'U'^\top=HU^\top
\end{gather}
are equivariant. The corresponding identities do not hold for general
$Q\in\mathrm O(q)$.
\end{proposition}
\begin{proof}
\begin{gather}
 (ZQ)^{\odot2}=Z^{\odot2}|Q|\quad(Q\text{ signed permutation}),\\
 (ZR)^{\odot2}\ne Z^{\odot2}R\quad\text{in general},\qquad
 (|ZR|+\epsilon)^{-1}\odot ZR\ne((|Z|+\epsilon)^{-1}\odot Z)R.
\end{gather}
\end{proof}
\begin{proposition}[Linear projector surrogate]\label{prop:tf-B}
Set $a_{K+1}:=a_\perp$ and
\begin{align}
 B(U)&:=\sum_{j=1}^Ka_j\Pi_j(U)+a_\perp\Pi_{K+1}(U)\\
 &=a_\perp I+\sum_{k=1}^K(a_k-a_{k+1})P_k(U),\\
 a_-I&\preceq B(U)\preceq a_+I,\quad
 a_-:=\min_j a_j>0,\quad a_+:=\max_j a_j,\\
 \Gamma^a&:=\sum_k|a_k-a_{k+1}|,\quad
 \|B(U)-B(V)\|_{\op}\le\Gamma^a d_P(U,V).
 \label{eq:tf-B-error}
\end{align}
In particular, $a_1=\cdots=a_K=a_\perp=a$ implies $B(U)=aI$.
\end{proposition}
\begin{proof}
\begin{align}
 x^\top B(U)x&=\sum_{j=1}^{K+1}a_j\|\Pi_jx\|^2,\quad
 \sum_j\|\Pi_jx\|^2=\|x\|^2,\\
 \|B(U)-B(V)\|_{\op}
 &\le\sum_k|a_k-a_{k+1}|\|P_k(U)-P_k(V)\|_F.
\end{align}
\end{proof}
For the linear surrogate, equal block coefficients remove all dependence on
the flag through $\Gamma^a=0$. Proposition~\ref{prop:tf-gauge} shows why the
corresponding Adam ablation has a narrower interpretation: entrywise
normalization still depends on the basis. The following decomposition keeps
moment lag and changing normalizers explicit when applying surrogate descent
results to the implemented update.

\begin{definition}[Exact predictable-surrogate decomposition]\label{def:rigid-error}
For any $\mathcal F_{t-1}$-measurable positive operator $B_t$ on
$\RR^{m\times d}$, define
\begin{gather}
 e_t^{\mathrm{alg}}:=D_t-B_tG_t,\qquad
 e_t:=e_t^{\mathrm{alg}}+\lambda W_t,\\
 W_{t+1}=W_t-\eta_t(B_tG_t+e_t).
\end{gather}
For a rowwise lagged diagonal surrogate, let
\begin{gather}
 h_{k,ij,t}^-:=(\sqrt{\widehat V_{k,ij,t-1}}+\epsilon)^{-1},\quad
 b_{i,t}^-:=\sum_k a_k U_{k,t}\diag(h_{k,i:,t}^-)U_{k,t}^\top
                     +a_\perp c_t^-\Pi_{K+1,t},\\
 (B_tG)_{i:}:=G_{i:}b_{i,t}^-,\qquad c_t^->0\text{ predictable}.
\end{gather}
Then, exactly,
\begin{align}
 e_t^{\mathrm{alg}}
 &=\sum_k a_k\left[
 (\widehat M_{k,t}-Z_{k,t})\odot h_{k,t}^-
 +\widehat M_{k,t}\odot(h_{k,t}-h_{k,t}^-)
 \right]U_{k,t}^\top\\
 &\quad+a_\perp(c_t-c_t^-)G_t\Pi_{K+1,t},\quad
 c_t:=(\|G_t\Pi_{K+1,t}\|_F/\sqrt{md}+\epsilon)^{-1},\\
 \|e_t^{\mathrm{alg}}\|_F
 &\le\sum_k|a_k|\left[
 \epsilon^{-1}\|\widehat M_{k,t}-Z_{k,t}\|_F+
 G_\star\|h_{k,t}-h_{k,t}^-\|_{\max}\right]
 +|a_\perp|G_\star|c_t-c_t^-|.
\end{align}
Set $h_{k,t}^-=\epsilon^{-1}\mathbf1$ when the lagged second-moment age is zero.
The moment-lag term includes the transport discrepancy of
Proposition~\ref{prop:rigid-transport}.
\end{definition}
\begin{lemma}[Clipping bias]\label{lem:rigid-clipping}
Let $\widetilde g$ be an unbiased gradient, $g=\widetilde g\min(1,c/\|\widetilde g\|)$,
$h=\EE\widetilde g$, and $\EE\|\widetilde g\|^p\le M_p$, $p>1$. Then
\begin{align}
 b_c&:=\EE g-h,\quad
 \|b_c\|\le\EE(\|\widetilde g\|-c)_+
 \le\frac{(p-1)^{p-1}}{p^p}\frac{M_p}{c^{p-1}}.
\end{align}
\end{lemma}
\begin{proof}
\begin{gather}
 \|g-\widetilde g\|=(\|\widetilde g\|-c)_+,\qquad
 \sup_{x>c}\frac{x-c}{x^p}
 =\frac{(p-1)^{p-1}}{p^pc^{p-1}},\quad x_{\max}=pc/(p-1).
\end{gather}
\end{proof}
\begin{proposition}[Second-order coordinate error for a horizontal QR step]
\label{prop:rigid-small-transport}
Let $Z\in\Hor_U$, $\overline U=\operatorname{qf}(U+Z)=U+Z+E$,
$e:=\|E\|_F<1$, and, blockwise, $C_j:=\overline U_j^\top U_j=O_jH_j$.
Then
\begin{gather}
 U_j^\top Z_j=0,\quad \|C_j-I\|_F\le e_j:=\|E_j\|_F,\quad
 H_j=U_j^\top(\overline U_jO_j),\\
 \|O_j-I\|_F\le\frac{2e_j}{2-e_j},\quad
 \|H_j-I\|_F\le e_j,\quad
 \|O_j^\top-U_j^\top(\overline U_jO_j)\|_F
                 \le\left(1+\frac2{2-e_j}\right)e_j,\\
 e\le(1+\|Z\|_{\op}^2)\|Z\|_F^2.
\end{gather}
Consequently the coordinate part of the transport discrepancy obeys
\begin{equation}
 \|M_j(O_j^\top-U_j^\top U_j^+)\|_F
 \le\|M_j\|_{\op}\left(1+\frac2{2-e}\right)
                    (1+\|Z\|_{\op}^2)\|Z\|_F^2.
\end{equation}
\end{proposition}
\begin{proof}
\begin{align}
 C_j&=I+Z_j^\top U_j+E_j^\top U_j=I+E_j^\top U_j,\\
 \sigma_{\min}(C_j)&\ge1-e_j,\qquad C_j^\top O_j=H_j,\\
 \|O_j-I\|_F&\le2\|C_j-I\|_F/(\sigma_{\min}(C_j)+1),\\
 \|C_j-I\|_F^2-\|H_j-I\|_F^2
 &=2\tr(H_j)-2\tr(O_jH_j)\ge0,\\
 \|O_j^\top-H_j\|_F&\le\|O_j-I\|_F+\|I-H_j\|_F.
\end{align}
\end{proof}

\subsection{Nonconvex convergence, PL rates, and coupled Lyapunov bounds}
We first analyze a predictable positive preconditioner with an explicit
discrepancy term. This separates the descent mechanism from the costs of
momentum, normalization, clipping, and transport. For STRATA the decomposition
is exact, but obtaining a vanishing error requires control of those costs.
The later analysis of the actual normalization field provides a second route
that retains its dependence on the current gradient directly.

\begin{assumption}[Smooth objective and predictable surrogate]\label{ass:rigid-descent}
For $F:\RR^D\to\RR$, $w_t\in\mathcal F_{t-1}$, and a self-adjoint $B_t$, assume
\begin{gather}
 F_{\inf}:=\inf_wF(w)>-\infty,\quad\EE|F(w_0)|<\infty,\\\quad
 \|\nabla F(w)-\nabla F(v)\|\le L\|w-v\|,\\
 h_t:=\nabla F(w_t),\quad g_t=h_t+b_t+\zeta_t,\quad
 \EE[\zeta_t\mid\mathcal F_{t-1}]=0,\quad
 \EE[\|\zeta_t\|^2\mid\mathcal F_{t-1}]\le\sigma_t^2,\\
 B_t\in\mathcal F_{t-1},\quad 0<mI\preceq B_t\preceq MI,\\
 w_{t+1}=w_t-\eta_t(B_tg_t+e_t),\quad
 \varepsilon_t:=e_t+B_tb_t.
\end{gather}
\end{assumption}
\begin{theorem}[Variance-resolved finite-time stationarity]\label{thm:tf-convergence}
Under Assumption~\ref{ass:rigid-descent}, let
$0<\eta_t\le1/(4LM)$, $S_T:=\sum_{t<T}\eta_t$, and
$\Delta_0:=\EE F(w_0)-F_{\inf}$. Define
\begin{equation}
 \Sigma_t:=\EE_t[\zeta_t\zeta_t^\top],\quad
 \mathsf v_t:=\tr(B_t^2\Sigma_t)=\EE_t\|B_t\zeta_t\|^2
                         \le M^2\sigma_t^2.
\end{equation}
Write $\|x\|_A^2:=x^\top Ax$ for $A\succ0$. Then
\begin{align}
 \EE F(w_{t+1})
 &\le\EE F(w_t)-\tfrac12\eta_t\EE\|h_t\|_{B_t}^2
       +L\eta_t^2\EE\mathsf v_t
       +\eta_t\EE\|\varepsilon_t\|_{B_t^{-1}}^2
       +L\eta_t^2\EE\|\varepsilon_t\|^2,\\
 \frac{\sum_{t<T}\eta_t\EE\|h_t\|^2}{S_T}
 &\le\frac{2\Delta_0}{mS_T}
 +\frac{2L\sum_{t<T}\eta_t^2\EE\mathsf v_t}{mS_T}
 +\frac{2\sum_{t<T}\left[\eta_t\EE\|\varepsilon_t\|_{B_t^{-1}}^2
                         +L\eta_t^2\EE\|\varepsilon_t\|^2\right]}{mS_T}.
 \label{eq:tf-stationarity}
\end{align}
For an independent output index $J$ with $\PP(J=t)=\eta_t/S_T$, the left side is
$\EE\|\nabla F(w_J)\|^2$.
\end{theorem}
\begin{proof}
Use the descent lemma~\citep{nesterov2004introductory} and the positive-operator
inequality $B_t^2\preceq MB_t$. Weighted Fenchel--Young bounds the defect
in the matching dual metric. With $\overline g_t:=h_t+\zeta_t$ and
$Q_t:=\langle h_t,B_th_t\rangle$,
\begin{align}
 F(w_{t+1})&\le F(w_t)-\eta_t\langle h_t,B_t\overline g_t+\varepsilon_t\rangle
                    +\tfrac L2\eta_t^2\|B_t\overline g_t+\varepsilon_t\|^2,\\
 \EE_t\langle h_t,B_t\overline g_t\rangle
 &=\langle h_t,B_th_t\rangle,\\
 |\langle h_t,\varepsilon_t\rangle|
 &\le\tfrac14 Q_t+\|\varepsilon_t\|_{B_t^{-1}}^2,\\
 \EE_t\|B_t\overline g_t+\varepsilon_t\|^2
 &\le2\|B_th_t\|^2+2\mathsf v_t+2\EE_t\|\varepsilon_t\|^2,\\
 \|B_th_t\|^2&\le M\langle h_t,B_th_t\rangle,\\
 -\eta_t(1-LM\eta_t)Q_t
                 +\tfrac14\eta_t Q_t
 &\le-\tfrac12\eta_t Q_t,\\
 \tfrac12\sum_{t<T}\eta_t\EE\|h_t\|_{B_t}^2
 &\le\EE F(w_0)-\EE F(w_T)
 +L\sum_{t<T}\eta_t^2\EE\mathsf v_t\\
 &\quad+\sum_{t<T}\left[\eta_t\EE\|\varepsilon_t\|_{B_t^{-1}}^2
                              +L\eta_t^2\EE\|\varepsilon_t\|^2\right].
\end{align}
Finally $\|h_t\|_{B_t}^2\ge m\|h_t\|^2$. The defect energy is kept directional;
replacing it by $m^{-1}\|\varepsilon_t\|^2$ is an optional later relaxation.
\end{proof}
\begin{corollary}[Sharp descent in the zero-discrepancy case]\label{cor:sharp-descent}
If $\varepsilon_t=0$ and $0<\eta_t<2/(LM)$, then
\begin{equation}
 \sum_{t<T}\eta_t(1-LM\eta_t/2)\EE\|h_t\|_{B_t}^2
 \le\Delta_0+\frac L2\sum_{t<T}\eta_t^2\EE\mathsf v_t.
\end{equation}
\end{corollary}
\begin{proof}
Conditional centering gives the exact second-moment identity
$\EE_t\|B_t(h_t+\zeta_t)\|^2=\|B_th_t\|^2+\mathsf v_t$.
Apply the descent lemma without a two-term norm relaxation and sum.
For $F(w)=L\|w\|^2/2$ and $B_t=MI$, every one-step inequality is an equality:
both the noise coefficient $L/2$ and the endpoint $2/(LM)$ are sharp for this
class. The positive-operator step also explains the larger step range in
Theorem~\ref{thm:tf-convergence}.
\end{proof}
The stationarity bound has three sources: the initial objective gap, accumulated
stochastic variance, and accumulated update discrepancy. Its left side concerns
a step-size-weighted output iterate. Thus a decreasing training loss alone
does not verify the conclusion, and persistent discrepancy can leave a
nonvanishing stationarity floor even as the stochastic term decreases.

\begin{corollary}[Rates and almost-sure stationarity along a subsequence]
If $\EE\mathsf v_t\le\overline v$, $\eta_t=\eta/\sqrt T\le1/(4LM)$, then
\begin{align}
 \EE\|\nabla F(w_J)\|^2
 &\le\frac{2\Delta_0}{m\eta\sqrt T}
       +\frac{2L\eta\overline v}{m\sqrt T}
       +\frac{2}{m}\left(\frac1m+\frac{L\eta}{\sqrt T}\right)
                        \frac1T\sum_{t<T}\EE\|\varepsilon_t\|^2.
\end{align}
If almost surely
\begin{gather}
 \sum_t\eta_t=\infty,\quad
 \sum_t\eta_t^2\mathsf v_t<\infty,\quad
 \sum_t(\eta_t/m+L\eta_t^2)\EE_t\|\varepsilon_t\|^2<\infty,
\end{gather}
then
\begin{equation}
 F(w_t)\text{ converges a.s.},\qquad
 \sum_t\eta_t\|h_t\|^2<\infty\text{ a.s.},\qquad
 \liminf_{t\to\infty}\|h_t\|=0\text{ a.s.}
\end{equation}
\end{corollary}
\begin{proof}
The finite-time assertion follows by substitution; one may always take
$\overline v=M^2\sigma^2$ when $\sigma_t^2\le\sigma^2$, and
$\|\varepsilon_t\|_{B_t^{-1}}^2\le m^{-1}\|\varepsilon_t\|^2$.
For the almost-sure assertion,
apply the nonnegative almost-supermartingale theorem \citep{robbins1971convergence}
to $F(w_t)-F_{\inf}$, the summable conditional error, and
$\tfrac m2\eta_t\|h_t\|^2$; an eventually positive lower bound on $\|h_t\|$
would contradict $\sum_t\eta_t=\infty$.
\end{proof}
\begin{theorem}[PL contraction and an explicit floor]\label{thm:tf-PL}
In addition, suppose
\begin{gather}
 \tfrac12\|\nabla F(w)\|^2\ge\mu(F(w)-F^\star),\quad\mu>0,\\
 0<\eta\le\min\{1/(4LM),1/(\mu m)\},\quad
 \EE\mathsf v_t\le\overline v,\quad\EE\|\varepsilon_t\|^2\le E^2.
\end{gather}
For $a_t:=\EE[F(w_t)-F^\star]$, $q:=1-\mu m\eta$,
\begin{align}
 a_{t+1}&\le qa_t+L\eta^2\overline v+(\eta/m+L\eta^2)E^2,\\
 a_T&\le q^Ta_0+
 \frac{L\eta\overline v+(m^{-1}+L\eta)E^2}{\mu m}(1-q^T).
\end{align}
\end{theorem}
\begin{proof}
\begin{gather}
 -\tfrac m2\eta\EE\|h_t\|^2\le-\mu m\eta a_t,\qquad
 \sum_{j=0}^{T-1}q^j=(1-q^T)/(\mu m\eta).
\end{gather}
\end{proof}
Under PL, the same descent estimate becomes a contraction of objective error.
The formula distinguishes a transient geometric term from the error floor:
reducing the learning rate suppresses the variance contribution, while a fixed
discrepancy budget $E^2$ can persist. The diminishing-step result specifies when
decaying discrepancy also permits the objective error to vanish.

\begin{lemma}[Diminishing-step PL recurrence]\label{lem:rigid-polynomial}
If $t_0>a>\max\{1,p\}$, $p>0$, and
\begin{equation}
 a_{t+1}\le\left(1-\frac a{t+t_0}\right)a_t+
                  \frac b{(t+t_0)^2}+\frac c{(t+t_0)^{1+p}},
\end{equation}
then
\begin{align}
 a_T&\le\left(\frac{t_0}{T+t_0}\right)^a a_0
 +\frac{b\,2^a}{(T+t_0)^a}\sum_{t<T}(t+t_0)^{a-2}
 +\frac{c\,2^a}{(T+t_0)^a}\sum_{t<T}(t+t_0)^{a-p-1}\\
 &=O((T+t_0)^{-a}+(T+t_0)^{-1}+(T+t_0)^{-p}).
\end{align}
Thus $\eta_t=\eta_0/(t+t_0)$, $\EE\|\varepsilon_t\|^2=O((t+t_0)^{-p})$,
and $a=\mu m\eta_0>\max\{1,p\}$ imply the displayed rate.
\end{lemma}
\begin{proof}
\begin{align}
 \prod_{j=s+1}^{T-1}\left(1-\frac a{j+t_0}\right)
 &\le\exp\left(-a\sum_{j=s+1}^{T-1}\frac1{j+t_0}\right)
 \le\left(\frac{s+1+t_0}{T+t_0}\right)^a,\\
 (s+1+t_0)^a&\le2^a(s+t_0)^a,\qquad
 \sum_{t<T}(t+t_0)^u=O((T+t_0)^{u+1})\quad(u>-1).
\end{align}
\end{proof}
\begin{theorem}[Coupled parameter--flag Lyapunov inequality]\label{thm:rigid-lyapunov}
Suppose nonnegative sequences $E_t,R_t,\Xi_t,\sigma_t^2$ satisfy
\begin{align}
 \EE F(w_{t+1})&\le\EE F(w_t)-a\eta\EE\|h_t\|^2
               +c\eta E_t+v\eta^2\sigma_t^2+r\eta R_t,\\
 E_{t+1}&\le\kappa E_t+b\eta^2\EE\|h_t\|^2+q\eta^2\sigma_t^2+s\Xi_t,
 \quad0\le\kappa<1.
\end{align}
Let $a,c>0$, $b\ge0$,
\begin{gather}
 \eta^2\le\frac{a(1-\kappa)}{4bc}\quad(bc=0:\ \text{no constraint}),\quad
 \Lambda:=\frac{2c\eta}{1-\kappa},\quad
 \mathcal L_t:=\EE F(w_t)-F_{\inf}+\Lambda E_t.
\end{gather}
Then
\begin{align}
 \mathcal L_{t+1}-\mathcal L_t
 &\le-\tfrac a2\eta\EE\|h_t\|^2-c\eta E_t
       +(v+\Lambda q)\eta^2\sigma_t^2+r\eta R_t+\Lambda s\Xi_t,\\
 \frac1T\sum_{t<T}\EE\|h_t\|^2
 &\le\frac{2\mathcal L_0}{a\eta T}
 +\frac{2(v+\Lambda q)\eta}{aT}\sum_{t<T}\sigma_t^2
 +\frac{2r}{aT}\sum_{t<T}R_t
 +\frac{4cs}{a(1-\kappa)T}\sum_{t<T}\Xi_t.
\end{align}
\end{theorem}
\begin{proof}
\begin{align}
 \mathcal L_{t+1}-\mathcal L_t
 &\le-(a\eta-\Lambda b\eta^2)\EE\|h_t\|^2
       +(c\eta-\Lambda(1-\kappa))E_t
       +(v+\Lambda q)\eta^2\sigma_t^2+r\eta R_t+\Lambda s\Xi_t,\\
 a\eta-\Lambda b\eta^2&\ge a\eta/2,\qquad
 c\eta-\Lambda(1-\kappa)=-c\eta.
\end{align}
Summation and $\mathcal L_T\ge0$ give the conclusion.
\end{proof}
The joint potential accounts for feedback between parameter motion and flag
tracking: each can introduce error into the other. Its weight $\Lambda$ is
chosen so that contraction of the flag error absorbs its contribution to
parameter descent. The following corollary then connects covariance estimation
and tracking to the stationarity bound through explicit intermediate constants.

\begin{corollary}[Covariance-to-stationarity error chain]\label{cor:tf-endtoend}
Suppose a linear block-isotropic surrogate satisfies
\begin{gather}
 \EE\|\varepsilon_t\|^2\le
   2G_\star^2(\Gamma^a)^2J_P^2 E_t+2R_t,\quad
 d_P(\psi(x),\psi(y))\le J_P\|x-y\|,\\
 E_{t+1}\le\kappa E_t+B_t.
\end{gather}
For constant $\eta\le1/(4LM)$ and $\sup_t\EE\mathsf v_t\le\overline v$,
\begin{align}
 \frac1T\sum_{t<T}\EE\|h_t\|^2
 &\le\frac{2\Delta_0}{m\eta T}+\frac{2L\eta\overline v}{m}\\
 &\quad+\frac4m(m^{-1}+L\eta)\left[
 \frac{G_\star^2(\Gamma^a)^2J_P^2}{T(1-\kappa)}
             \left(E_0+\sum_{t<T-1}B_t\right)+\frac1T\sum_{t<T}R_t\right].
\end{align}
For the implemented coordinatewise optimizer, $R_t$ additionally contains the
normalization, gauge, momentum, clipping, and decay terms of
Definitions~\ref{def:rigid-map} and~\ref{def:rigid-error} and
Lemma~\ref{lem:rigid-clipping}.
\end{corollary}
\begin{proof}
\begin{gather}
 \|(B(U_t)-B(U_t^\star))g_t\|^2
 \le G_\star^2(\Gamma^a)^2J_P^2\|x_t-x_t^\star\|^2,\quad
 \sum_{t<T}E_t\le(E_0+\sum_{t<T-1}B_t)/(1-\kappa).
\end{gather}
\end{proof}
\subsection{A virtual-iterate theorem for matrix momentum}
The virtual iterate subtracts the parameter displacement stored in momentum.
Its increment resembles a preconditioned gradient step, with a remainder that
records changing coefficients and transport. This algebraic reformulation
allows the proof to retain the whole momentum history through variation
budgets, rather than discarding it as an unspecified per-step error.

\begin{assumption}[Predictable momentum system]\label{ass:rigid-momentum}
Let $\mathcal B_t,B_t\in\mathcal F_{t-1}$, with
$\|\mathcal B_t\|_{\op}\le\beta<1$, $mI\preceq B_t\preceq MI$, and
\begin{gather}
 m_t=\mathcal B_t m_{t-1}+(I-\mathcal B_t)g_t+u_t,\quad
 w_{t+1}=w_t-\eta_tB_tm_t,\quad m_{-1}=0,\\
 \EE_tg_t=\nabla F(w_t),\quad\EE_t\|g_t\|^2\le G_2^2,\quad
 \|m_t\|\le G_m,\quad\|\nabla F(z_t)\|\le G_z,\\
 C_t:=\eta_tB_t(I-\mathcal B_t)^{-1}\mathcal B_t,\quad C_{-1}:=0,\quad
 z_t:=w_t-C_{t-1}m_{t-1}.
\end{gather}
\end{assumption}
\begin{lemma}[Exact virtual-iterate identity]\label{lem:rigid-virtual}
Under Assumption~\ref{ass:rigid-momentum},
\begin{gather}
 z_{t+1}=z_t-\eta_tB_tg_t+r_t,\\
 r_t:=(C_{t-1}-C_t)m_{t-1}-\eta_tB_t(I-\mathcal B_t)^{-1}u_t,\\
 \|r_t\|\le G_m\|C_t-C_{t-1}\|_{\op}
                      +\frac{\eta_tM}{1-\beta}\|u_t\|,\qquad
 \|w_t-z_t\|\le\|C_{t-1}\|_{\op}G_m.
\end{gather}
\end{lemma}
\begin{proof}
\begin{align}
 \eta_tB_t+C_t&=\eta_tB_t(I-\mathcal B_t)^{-1},\\
 (\eta_tB_t+C_t)m_t
 &=C_tm_{t-1}+\eta_tB_tg_t
                       +\eta_tB_t(I-\mathcal B_t)^{-1}u_t,\\
 z_{t+1}-z_t&=-(\eta_tB_t+C_t)m_t+C_{t-1}m_{t-1}.
\end{align}
\end{proof}
\begin{theorem}[Stationarity with explicit variation and transport budgets]
\label{thm:rigid-momentum}
Under Assumptions~\ref{ass:rigid-descent} (smoothness and lower bound only)
and~\ref{ass:rigid-momentum},
\begin{align}
 \frac m2\sum_{t<T}\eta_t\EE\|\nabla F(w_t)\|^2
 &\le F(w_0)-F_{\inf}+LM^2G_2^2\sum_{t<T}\eta_t^2\\
 &\quad+\frac{M^2L^2G_m^2}{2m}\sum_{t<T}\eta_t\|C_{t-1}\|_{\op}^2
       +G_z\sum_{t<T}\EE\|r_t\|+L\sum_{t<T}\EE\|r_t\|^2.
\end{align}
For constant scalar momentum $\mathcal B_t=\beta I$ and constant
$\eta=\eta_0/\sqrt T$, define
$\mathcal V_T:=\sum_{t<T}\|B_t-B_{t-1}\|_{\op}$, with $B_{-1}=0$. Then
\begin{align}
 \sum_{t<T}\EE\|r_t\|
 &\le\frac{\eta\beta G_m}{1-\beta}\EE\mathcal V_T
                  +\frac{\eta M}{1-\beta}\sum_{t<T}\EE\|u_t\|,\\
 \EE\mathcal V_T+\sum_{t<T}\EE\|u_t\|=O(\sqrt T),\quad
 \sum_{t<T}\EE\|r_t\|^2=O(1)
 &\Longrightarrow\frac1T\sum_{t<T}\EE\|\nabla F(w_t)\|^2=O(T^{-1/2}).
\end{align}
\end{theorem}
\begin{proof}
Let $h_t=\nabla F(w_t)$, $\delta_t=\nabla F(z_t)-h_t$. Then
\begin{align}
 \|\delta_t\|&\le L\|C_{t-1}\|_{\op}G_m,\\
 -\EE_t\langle\nabla F(z_t),B_tg_t\rangle
 &=-\langle h_t+\delta_t,B_th_t\rangle\\
 &\le-\tfrac m2\|h_t\|^2+\tfrac{M^2}{2m}\|\delta_t\|^2,\\
 F(z_{t+1})&\le F(z_t)-\eta_t\langle\nabla F(z_t),B_tg_t\rangle
 +\langle\nabla F(z_t),r_t\rangle
 +\tfrac L2\|-\eta_tB_tg_t+r_t\|^2,\\
 \EE_tF(z_{t+1})&\le F(z_t)-\tfrac m2\eta_t\|h_t\|^2
 +\frac{\eta_tM^2L^2G_m^2}{2m}\|C_{t-1}\|_{\op}^2
 \notag\\
 &\quad+G_z\EE_t\|r_t\|+LM^2\eta_t^2G_2^2+L\EE_t\|r_t\|^2.
\end{align}
Summation gives the first assertion. Also
$\|C_t\|_{\op}\le\eta M\beta/(1-\beta)$ and
$\sum_t\eta\|C_{t-1}\|_{\op}^2=O(T\eta^3)$.
\end{proof}

\subsection{Full-state replace-one stability and generalization}
Replacing one training example perturbs more than the current parameter
gradient. It also changes covariance sketches, subsequent flags, and both
moment states, which can feed back into later parameter updates. We therefore
couple two runs with shared randomness and track five distances. The
nonnegative transition matrix records how each component can amplify the
others over the remaining training horizon.

\begin{definition}[Neighboring samples and synchronized coupling]
\begin{gather}
 S=(z_1,\ldots,z_n)\sim\mathcal D^n,\quad z_i'\sim\mathcal D,\quad
 S^{(i)}=(z_1,\ldots,z_{i-1},z_i',z_{i+1},\ldots,z_n),\\
 F_S(w):=\frac1n\sum_{i=1}^n\ell(w,z_i),\quad
 F_\mathcal D(w):=\EE_{z\sim\mathcal D}\ell(w,z),\\
 A_T(S;\vartheta)=w_T,\quad A_T(S^{(i)};\vartheta)=w_T',\quad
 \vartheta=(U_0,\Omega_{1:K},\text{batch indices, augmentations, masks}),\\
 I_t:=\mathbf1_{i\in\mathcal I_t},\quad |\mathcal I_t|=b_t,\quad
 \PP(I_t=1)=b_t/n.
\end{gather}
The initial state, probes, schedule, counters, and random seed are shared.
The marginal sampling condition permits fresh uniform subsets and uniform random
permutations; conditional independence from past states is not required below.
\end{definition}
\begin{assumption}[Uniform regularity for a coupled pair]\label{ass:rigid-stability}
On the trajectories under consideration, assume deterministic envelopes
\begin{gather}
 |\ell(w,z)-\ell(v,z)|\le L_0\|w-v\|,\quad
 \|\nabla\ell(w,z)\|\le G_\star,\quad
 \|\nabla\ell(w,z)-\nabla\ell(v,z)\|\le L_g\|w-v\|,\\
 \Omega_k^\top\Omega_k=I,\quad
 \sym(\Omega_k^\top Y_k)+\gamma I\succeq\gamma I,\quad\|Y_k\|_{\op}\le Y_\star,\\
 \|\widetilde C_k\|_{\op}\le C_\star,\quad
 \inf_{s\in[0,1]}\sigma_{\min}(X+s(X'-X))\ge a_Q>0,\\
 X:=U+\xi\operatorname{grad}\Phi(U;\widetilde C),\quad
 \min\{\sigma_{\min}(\overline U_k^\top U_k),
       \sigma_{\min}(\overline U_k'^\top U_k')\}\ge a_O>0,\\
 0\le\eta_t\lambda\le1,\quad a_k\ge0,\quad a_\perp\ge0.
\end{gather}
The gradient passed to the optimizer may additionally undergo the same
nonexpansive norm clipping in both runs.
\end{assumption}
\begin{definition}[Five distances and deterministic constants]
\begin{gather}
 u_t:=\|w_t-w_t'\|,\quad v_t:=\|U_t-U_t'\|_F,\quad
 y_t:=\sum_k\omega_k\|Y_{k,t-1}-Y_{k,t-1}'\|_F,\\
 m_t:=\sum_k\|M_{k,t-1}^+-M_{k,t-1}^{+\prime}\|_F,\quad
 s_t:=\sum_k\|\sqrt{V_{k,t-1}^+}-\sqrt{V_{k,t-1}^{+\prime}}\|_F,\\
 z_t:=(u_t,v_t,y_t,m_t,s_t)^\top\in\RR_+^5,\\
 \bar\rho:=\max_k\rho_k,\quad
 \bar\beta_j:=\max_k\beta_{jk},\quad\underline\beta_j:=\min_k\beta_{jk},\\
 c_Y:=2G_\star\sum_k\omega_k(1-\rho_k),\quad
 c_M:=\sqrt K(1-\underline\beta_1),\quad
 c_S:=\sqrt K\sqrt{1-\underline\beta_2},\\
 H_\star:=2\sqrt R C_\star\sum_k\omega_k,\quad
 L_U:=2C_\star\sum_k\omega_k+4H_\star,\quad
 L_Y:=2\sqrt R\max_k L_{N,k},\\
 L_Q:=\sqrt2/a_Q,\quad L_O:=1/a_O,\quad c_O:=\sqrt K L_O,\\
 A_M:=((1-\bar\beta_1)\epsilon)^{-1},\quad
 A_S:=\frac{G_\star}{(1-\bar\beta_1)\epsilon^2\sqrt{1-\bar\beta_2}},\quad
 H_0:=G_\star/\epsilon,\quad a_\star:=\max_k a_k.
\end{gather}
For multiple matrices and dense fallback parameters, concatenate their distances;
the same construction gives a larger nonnegative matrix, or a five-state bound
using valid aggregate envelopes.
\end{definition}
\begin{lemma}[Gradient, covariance, and horizontal-gradient perturbations]
\label{lem:rigid-state-primitives}
Under Assumption~\ref{ass:rigid-stability},
\begin{align}
 \delta G_t&:=\|G_t-G_t'\|_F\le L_gu_t+(2G_\star/b_t)I_t,\\
 \|G_t^\top G_t-G_t'^\top G_t'\|_F&\le2G_\star\delta G_t,\\
 y_t^-&:=\sum_k\omega_k\|Y_{k,t}^--Y_{k,t}^{-\prime}\|_F
             \le\bar\rho y_t+c_Y\delta G_t,\\
 \|\operatorname{grad}\Phi(U;\widetilde C)
        -\operatorname{grad}\Phi(U';\widetilde C')\|_F
 &\le L_U\|U-U'\|_F+L_Y\sum_k\omega_k\|Y_k-Y_k'\|_F,\\
 \sum_k\|Z_k-Z_k'\|_F&\le\sqrt K(\delta G_t+G_\star v_t).
\end{align}
\end{lemma}
\begin{proof}
\begin{align}
 G_t-G_t'&=\frac1{b_t}\sum_{j\in\mathcal I_t}
       [\nabla\ell(w_t,z_j)-\nabla\ell(w_t',z_j)]
       \notag\\
 &\quad+\frac{I_t}{b_t}[\nabla\ell(w_t',z_i)-\nabla\ell(w_t',z_i')],\\
 G^\top G-G'^\top G'&=(G-G')^\top G+G'^\top(G-G'),\\
 \|\mathsf H_U(H)-\mathsf H_{U'}(H)\|_F
 &\le4\|H\|_F\|U-U'\|_F,\quad\|\mathsf H_U\|_{F\to F}=1,\\
 \|H(U;C)-H(U';C')\|_F
 &\le2C_\star\sum_k\omega_k\|U-U'\|_F
       +2\sum_k\omega_k\sqrt{r_k}\|C_k-C_k'\|_{\op},\\
 \sum_k\|(G-G')U_k\|_F&\le\sqrt K\|(G-G')U\|_F\le\sqrt K\delta G_t,\\
 \sum_k\|G'(U_k-U_k')\|_F&\le\sqrt K G_\star\|U-U'\|_F.
\end{align}
\end{proof}
\begin{definition}[An explicit nonnegative transition matrix]\label{def:rigid-transition}
Let $e_w,e_U,e_Y,e_M,e_S$ be the standard row vectors in $\RR^5$.
The following rows and scalars define $\mathcal M_t\in\RR_+^{5\times5}$ and
$d_t\in\RR_+^5$ without implicit constants:
\begin{align}
 r_g&:=L_g e_w,& d_g&:=2G_\star/b_t,\\
 r_y&:=\bar\rho e_Y+c_Yr_g,& d_y&:=c_Yd_g,\\
 r_h&:=L_Q[(1+\xi L_U)e_U+\xi L_Yr_y],&d_h&:=L_Q\xi L_Yd_y,\\
 r_o&:=\chi_t c_O(r_h+e_U),&d_o&:=\chi_t c_Od_h,\\
 r_U&:=(1-\chi_t)e_U+\chi_t r_h+r_o,&d_U&:=\chi_t d_h+d_o,\\
 r_m^-&:=\bar\beta_1 e_M+c_M(r_g+G_\star e_U),&d_m^-&:=c_Md_g,\\
 r_s^-&:=\sqrt{\bar\beta_2}e_S+c_S(r_g+G_\star e_U),&d_s^-&:=c_Sd_g,\\
 r_M&:=r_m^-+G_\star r_o,&d_M&:=d_m^-+G_\star d_o,\\
 r_S&:=(1-\chi_t)r_s^-,&d_S&:=(1-\chi_t)d_s^-,\\
 \end{align}
\begin{align}
 r_D&:=a_\star(A_Mr_m^-+A_Sr_s^-+H_0\sqrt K e_U)\notag\\
 &\quad+(a_\perp/\epsilon)(r_g+2G_\star e_U),\\
 d_D&:=a_\star(A_Md_m^-+A_Sd_s^-)+(a_\perp/\epsilon)d_g,\\
 r_w&:=(1-\eta_t\lambda)e_w+\eta_t r_D,\\
 d_w&:=\eta_t d_D.
\end{align}
\begin{gather}
 \mathcal M_t:=\begin{pmatrix}r_w\\r_U\\r_y\\r_M\\r_S\end{pmatrix},\quad
 d_t:=(d_w,d_U,d_y,d_M,d_S)^\top.
\end{gather}
\end{definition}
\begin{theorem}[Full-state stability with all feedback paths]\label{thm:tf-stability}
Under Assumption~\ref{ass:rigid-stability}, Definition~\ref{def:rigid-map} satisfies,
componentwise,
\begin{align}
 z_{t+1}&\preceq\mathcal M_tz_t+d_tI_t,\qquad z_0=0,\\
 \EE z_T&\preceq\sum_{t=0}^{T-1}
       (\mathcal M_{T-1}\cdots\mathcal M_{t+1})d_t\frac{b_t}{n},\\
 \epsilon_T^{\mathrm{stab}}
 &:=\sup_{S\sim S',z}\EE_\vartheta
       |\ell(A_T(S;\vartheta),z)-\ell(A_T(S';\vartheta),z)|\\
 &\le\frac{L_0}{n}e_w\sum_{t=0}^{T-1}b_t
             (\mathcal M_{T-1}\cdots\mathcal M_{t+1})d_t.
 \label{eq:rigid-stability-product}
\end{align}
Empty products equal $I_5$. The supremum is over pairs satisfying the stated
uniform envelopes; a uniform result over all neighboring samples requires those
envelopes uniformly over the sample domain.
\end{theorem}
\begin{proof}
At a refresh, Lemmas~\ref{lem:rigid-qr-lip} and~\ref{lem:rigid-polar} imply
\begin{align}
 \|\overline U-\overline U'\|_F&\le r_hz_t+d_hI_t,\\
 \sum_k\|O_k-O_k'\|_F
 &\le\sqrt K L_O(\|\overline U-\overline U'\|_F+\|U-U'\|_F)
 \le r_oz_t+d_oI_t,\\
 \|\overline UO-\overline U'O'\|_F
 &\le\|\overline U-\overline U'\|_F+\sum_k\|O_k-O_k'\|_F,\\
 \sum_k\|M_k^-O_k^\top-M_k'^{-}O_k'^\top\|_F
 &\le\sum_k\|M_k^--M_k'^{-}\|_F+G_\star\sum_k\|O_k-O_k'\|_F.
\end{align}
The root-state inequality and the synchronized reset give
\begin{align}
 \sum_k\|M_k^--M_k'^{-}\|_F&\le r_m^-z_t+d_m^-I_t,\\
 \sum_k\|\sqrt{V_k^-}-\sqrt{V_k'^{-}}\|_F&\le r_s^-z_t+d_s^-I_t,\\
 \sum_k\|\sqrt{V_k^+}-\sqrt{V_k'^{+}}\|_F
 &=(1-\chi_t)\sum_k\|\sqrt{V_k^-}-\sqrt{V_k'^{-}}\|_F.
\end{align}
For the parameter direction,
\begin{align}
 \|HU^\top-H'U'^\top\|_F
 &\le\|H-H'\|_F+\|H'\|_{\op}\|U-U'\|_F,\\
 \|G(I-UU^\top)-G'(I-U'U'^\top)\|_F
 &\le\delta G_t+2G_\star v_t,\\
 \|D_t-D_t'\|_F&\le r_Dz_t+d_DI_t,\\
 u_{t+1}&\le(1-\eta_t\lambda)u_t+\eta_t\|D_t-D_t'\|_F.
\end{align}
Every entry of $\mathcal M_t,d_t$ is deterministic and nonnegative. Consequently,
\begin{align}
 \EE z_{t+1}&\preceq\mathcal M_t\EE z_t+d_t\EE I_t
 =\mathcal M_t\EE z_t+d_tb_t/n,\\
 \EE|\ell(w_T,z)-\ell(w_T',z)|&\le L_0 e_w\EE z_T.
\end{align}
\end{proof}
The matrix product retains every feedback path after the changed example is
sampled. A scalar bound follows if a weighted combination of state distances
has controlled growth. Such a weight is an additional property of the coupled
dynamics; bounding each one-step map by smoothness alone need not produce
contraction.

\begin{corollary}[Weighted small-gain and contraction bounds]\label{cor:rigid-smallgain}
Suppose a fixed row $a\in\RR_{>0}^5$, $a_w\ge1$, satisfies
\begin{gather}
 a\mathcal M_t\preceq(1+\nu\eta_t)a,\qquad
 b_t a d_t\le c\eta_t.
\end{gather}
Then
\begin{align}
 \epsilon_T^{\mathrm{stab}}
 &\le\frac{L_0c}{n}\sum_{t<T}\eta_t\prod_{s=t+1}^{T-1}(1+\nu\eta_s)\\
 &\le\begin{cases}
 \dfrac{L_0c}{n\nu}(e^{\nu S_T}-1),&\nu>0,\\[2pt]
 \dfrac{L_0c}{n}S_T,&\nu=0.
 \end{cases}
\end{align}
If instead $a\mathcal M_t\preceq(1-\nu\eta_t)a$, $0<\nu\eta_t\le1$, then
\begin{equation}
 \epsilon_T^{\mathrm{stab}}\le\frac{L_0c}{n\nu}
             \left[1-\prod_{t<T}(1-\nu\eta_t)\right]\le\frac{L_0c}{n\nu}.
\end{equation}
\end{corollary}
\begin{proof}
\begin{gather}
 a\EE z_{t+1}\le(1\pm\nu\eta_t)a\EE z_t+c\eta_t/n,\qquad
 \sum_{t<T}\nu\eta_t\prod_{s>t}(1\pm\nu\eta_s)
 =\begin{cases}\prod_t(1+\nu\eta_t)-1,&+,\\1-\prod_t(1-\nu\eta_t),&-.
 \end{cases}
\end{gather}
\end{proof}
\begin{proposition}[Regularity failure probability]\label{prop:rigid-exit}
If $0\le\ell\le B_\ell$, and the coupled regularity conditions fail before $T$
with probability at most $p_T$, then the same matrix-product upper bound, denoted
$\epsilon_{T,\mathrm{reg}}$, gives
\begin{equation}
 \epsilon_T^{\mathrm{stab}}\le\epsilon_{T,\mathrm{reg}}+B_\ell p_T.
\end{equation}
\end{proposition}
\begin{proof}
Let $\tau_\mathrm{exit}$ be the first failure time. Nonnegative killed distances
$\widetilde z_t:=z_t\mathbf1_{\tau_\mathrm{exit}>t}$ satisfy
\begin{gather}
 \widetilde z_{t+1}\preceq\mathcal M_t\widetilde z_t+d_tI_t,\\
 \EE|\ell(w_T,z)-\ell(w_T',z)|
 \le L_0e_w\EE\widetilde z_T+B_\ell\PP(\tau_\mathrm{exit}\le T).
\end{gather}
\end{proof}
\begin{theorem}[Expected generalization and concentration]\label{thm:rigid-generalization}
For iid $S$, independent algorithm randomness, and uniform replace-one stability
$\epsilon_T^{\mathrm{stab}}=\epsilon_T$,
\begin{equation}\label{eq:tf-gen}
 \left|\EE_{S,\vartheta}\left[F_\mathcal D(A_T(S;\vartheta))
                          -F_S(A_T(S;\vartheta))\right]\right|\le\epsilon_T.
\end{equation}
If also $0\le\ell\le B_\ell$, define
$f(S):=\EE_\vartheta[F_\mathcal D(A_T(S))-F_S(A_T(S))]$. Then, with probability
at least $1-\delta$ over $S$,
\begin{equation}
 |f(S)|\le\epsilon_T+(2n\epsilon_T+B_\ell)
                         \sqrt{\frac{\log(2/\delta)}{2n}}.
\end{equation}
For a bound on a single randomized run, the latter assertion requires the same
uniform stability constant for every fixed seed $\vartheta$; it then holds over
$(S,\vartheta)$.
\end{theorem}
\begin{proof}
Exchangeability and independence imply
\begin{align}
 \EE_{S,\vartheta}F_\mathcal D(A_T(S))
 &=\frac1n\sum_i\EE_{S,z_i',\vartheta}\ell(A_T(S),z_i'),\\
 \EE_{S,z_i',\vartheta}\ell(A_T(S^{(i)}),z_i')
 &=\EE_{S,\vartheta}\ell(A_T(S),z_i),\\
 \EE(F_\mathcal D-F_S)
 &=\frac1n\sum_i\EE[\ell(A_T(S),z_i')-\ell(A_T(S^{(i)}),z_i')].
\end{align}
For neighboring $S,S'$,
\begin{align}
 |f(S)-f(S')|
 &\le\EE_\vartheta|F_\mathcal D(A_T(S))-F_\mathcal D(A_T(S'))|\\
 &\quad+\frac1n\sum_{j\ne i}\EE_\vartheta
             |\ell(A_T(S),z_j)-\ell(A_T(S'),z_j)|+B_\ell/n\\
 &\le(2-1/n)\epsilon_T+B_\ell/n\le2\epsilon_T+B_\ell/n.
\end{align}
The bounded-differences inequality \citep{mcdiarmid1989bounded} yields
\begin{gather}
 \PP(|f(S)-\EE f(S)|\ge r)
 \le2\exp\left(-\frac{2r^2}{n(2\epsilon_T+B_\ell/n)^2}\right),\quad
 |\EE f(S)|\le\epsilon_T.
\end{gather}
\end{proof}
\begin{corollary}[Bousquet--Klochkov--Zhivotovskiy concentration]\label{cor:sharp-stability-tail}
Under the bounded-loss and uniform expected-stability hypotheses above, there
is a numerical constant $C$ such that, for $n\ge2$, with probability at least
$1-\delta$ over $S$,
\begin{equation}
 |f(S)|\le C\left[
 \epsilon_T\log(2n)\log(2/\delta)
       +B_\ell\sqrt{\frac{\log(2/\delta)}n}\right].
\end{equation}
If the stability bound holds for every fixed random seed with the same
constant, this bound also applies to a single run over $(S,\vartheta)$.
\end{corollary}
\begin{proof}
Define the deterministic, sample-dependent loss
$\overline\ell_S(z):=\EE_\vartheta\ell(A_T(S;\vartheta),z)$.
It takes values in $[0,B_\ell]$ and satisfies
\begin{equation}
 |\overline\ell_S(z)-\overline\ell_{S'}(z)|
 \le\EE_\vartheta|\ell(A_T(S;\vartheta),z)-\ell(A_T(S';\vartheta),z)|
 \le\epsilon_T.
\end{equation}
Apply Corollary 8 of \citet{bousquet2020sharper} to this uniformly stable
loss-valued algorithm. Its generalization gap is exactly $f(S)$. For the
single-run assertion condition on the seed and integrate the failure
probability. This improves the stability dependence from
$\epsilon_T\sqrt n$ in the elementary bounded-differences bound to
$\epsilon_T\log n$, up to confidence factors and a universal constant;
we make no claim that the numerical constant $C$ is optimal.
\end{proof}

\subsection{Effective noise dimension, excess risk, and horizon dependence}
The next results separate directional noise reduction from generalization and
from state storage. The trace identity measures how a predictable
preconditioner transforms gradient noise; the excess-risk bound also needs
optimization and stability estimates. Compressing optimizer state leaves the
parameter domain unchanged, as Proposition~\ref{prop:rigid-domain} makes explicit.

\begin{proposition}[Exact directional-noise identities]\label{prop:rigid-noise}
For predictable symmetric $B_t$ and conditional noise covariance
$\Sigma_t:=\EE_t[\zeta_t\zeta_t^\top]$,
\begin{gather}
 \EE_t\|B_t\zeta_t\|^2=\tr(B_t^2\Sigma_t),\quad
 \lambda_{\max}(\Sigma_t)\tr(B_t^2)\ge\tr(B_t^2\Sigma_t),\\
 B=\sum_{j=1}^{K+1}a_j\Pi_j\quad\Longrightarrow\quad
 \tr(B^2)=\sum_{j=1}^{K+1}q_ja_j^2.
\end{gather}
For $m$ rows with independent isotropic noise variance $\sigma^2$ per coordinate,
\begin{equation}
 \EE\|\zeta B\|_F^2=m\sigma^2
       \left(\sum_{j=1}^Kq_ja_j^2+(d-R)a_\perp^2\right).
\end{equation}
\end{proposition}
\begin{proof}
\begin{align}
 \EE_t\|B_t\zeta_t\|^2
 &=\EE_t\tr(B_t\zeta_t\zeta_t^\top B_t)=\tr(B_t^2\Sigma_t),\\
 B^2&=\sum_ja_j^2\Pi_j,\quad\tr\Pi_j=q_j.
\end{align}
\end{proof}
\begin{corollary}[Population excess risk under empirical PL]\label{cor:rigid-risk}
Suppose every $F_S$ is $\mu$-PL, a random output $A_J$ satisfies
$\EE\|\nabla F_S(A_J)\|^2\le R_T$, and its uniform expected stability is
$\epsilon_T$. Then for every fixed population minimizer $w_\mathcal D^\star$,
\begin{equation}
 \EE[F_\mathcal D(A_J)-F_\mathcal D(w_\mathcal D^\star)]
 \le R_T/(2\mu)+\epsilon_T.
\end{equation}
If, under the additional hypotheses of the stationarity and small-gain results,
\begin{gather}
 R_T/(2\mu)\le A/\sqrt T+R_{\mathrm{disc}},\quad
 \epsilon_T\le B\sqrt T/n,
\end{gather}
then
\begin{gather}
 \EE[F_\mathcal D(A_J)-F_\mathcal D(w_\mathcal D^\star)]
 \le A/\sqrt T+B\sqrt T/n+R_{\mathrm{disc}},\\
 T_\mathrm{cont}^\star=An/B,\quad
 \inf_{T>0}\{A/\sqrt T+B\sqrt T/n\}=2\sqrt{AB/n}.
\end{gather}
The two assumed rates must hold with constants uniform in the selected horizon.
\end{corollary}
\begin{proof}
\begin{align}
 F_S(A_J)-\inf_wF_S(w)&\le\|\nabla F_S(A_J)\|^2/(2\mu),\\
 \EE\inf_wF_S(w)&\le\EE F_S(w_\mathcal D^\star)
 =F_\mathcal D(w_\mathcal D^\star),\\
 \EE[F_\mathcal D(A_J)-F_\mathcal D(w_\mathcal D^\star)]
 &\le\EE[F_\mathcal D(A_J)-F_S(A_J)]
       +\EE[F_S(A_J)-\inf_wF_S(w)],\\
 \frac{\mathrm d}{\mathrm dT}(AT^{-1/2}+BT^{1/2}/n)
 &=-\tfrac A2T^{-3/2}+\tfrac B{2n}T^{-1/2}.
\end{align}
\end{proof}
The horizon balance expresses a conditional early-stopping tradeoff: more
updates reduce optimization error while the stability term can increase.
Its minimizing horizon is meaningful only when both assumed rates hold
uniformly over the horizons being compared. The movement bound below remains
available with fewer state-regularity requirements, though it can be much looser.

\begin{proposition}[A movement-only bound]\label{prop:tf-movement}
For a deterministic common initialization $w_0$, $L_0$-Lipschitz losses, and
$w_{t+1}=w_t-\eta_td_t$,
\begin{equation}
 \left|\EE[F_\mathcal D(w_T)-F_S(w_T)]\right|
 \le2L_0\sum_{t<T}\eta_t\EE\|d_t\|.
\end{equation}
\end{proposition}
\begin{proof}
\begin{gather}
 \EE[F_\mathcal D(w_0)-F_S(w_0)]=0,\quad
 |(F_\mathcal D-F_S)(w_T)-(F_\mathcal D-F_S)(w_0)|\le2L_0\|w_T-w_0\|,\\
 \|w_T-w_0\|\le\sum_{t<T}\eta_t\|d_t\|.
\end{gather}
\end{proof}
\begin{proposition}[State dimension versus parameter domain]\label{prop:rigid-domain}
For $m\ge d$, the stored floating-point matrix-state counts are
\begin{gather}
 N_{\mathrm{STRATA}}=2mR+dR+2d\sum_{k=1}^K\min(d,r_k+p),\qquad
 N_{\mathrm{Adam}}=2md,\\
 \dim\mathfrak F=dR-\tfrac12(R^2+\sum_kq_k^2),\qquad
 \mathcal W_{\mathrm{STRATA}}=\mathcal W_{\mathrm{Adam}}=\RR^{m\times d},\\
 U\in\St(d,R)\quad\not\Longrightarrow\quad \operatorname{rank}(W)\le R.
\end{gather}
Fixed probes are included in $N_{\mathrm{STRATA}}$; temporary workspaces,
parameters, gradients, scalar counters, and fallback states are excluded.
\end{proposition}
\begin{proof}
\begin{gather}
 N_M+N_V=2m\sum_kq_k=2mR,\quad N_U=dR,\quad
 N_Y=N_\Omega=d\sum_k\min(d,r_k+p),\\
 W_0=\begin{pmatrix}I_d\\0\end{pmatrix}\quad\Longrightarrow\quad
 \operatorname{rank}(W_0)=d>R,\qquad U_0^\top U_0=I_R.
\end{gather}
No projection of $W$ onto a rank-constrained set occurs in
Definition~\ref{def:rigid-map}.
\end{proof}
\begin{corollary}[Explicit local regularity radii]\label{cor:rigid-regularity-radius}
At a refresh, let $Z=\xi\operatorname{grad}\Phi(U;\widetilde C)$ and
$Z'=\xi\operatorname{grad}\Phi(U';\widetilde C')$. If
\begin{gather}
 \max(\|Z\|_F,\|Z'\|_F)\le z_\star,\quad
 e_\star:=(1+z_\star^2)z_\star^2<1,\quad
 \|(U+Z)-(U'+Z')\|_{\op}\le r_X<1,
\end{gather}
then the QR and Procrustes regularity assumptions hold with
\begin{equation}
 a_Q=1-r_X,\qquad a_O=1-e_\star.
\end{equation}
If $\|\mathrm D\psi\|\le J_\psi$ on a convex chart domain, then also
\begin{equation}
 d_P(\psi(x),\psi(y))\le\sqrt{2K}J_\psi\|x-y\|.
\end{equation}
\end{corollary}
\begin{proof}
\begin{align}
 \sigma_{\min}(U+Z)^2
 &=\lambda_{\min}(I+Z^\top Z)\ge1,\\
 \sigma_{\min}(X+s(X'-X))&\ge\sigma_{\min}(X)-s\|X'-X\|_{\op}\ge1-r_X,\\
 \sigma_{\min}(\overline U_j^\top U_j)
 &\ge1-\|\overline U_j-U_j-Z_j\|_F\ge1-e_\star,\\
 \left\|\frac{\mathrm d}{\mathrm ds}(P_1,\ldots,P_K)
                     (\psi(y+s(x-y)))\right\|_F
 &\le\sqrt{2K}J_\psi\|x-y\|.
\end{align}
\end{proof}

\subsection{Invariant-block quadratic rates and an exact small-gain criterion}
An invariant-block quadratic isolates a setting where differentiated scalar
weights can improve conditioning exactly. The result describes the benefit
available when the flag aligns with Hessian invariant subspaces. Learning a
flag from gradient covariance does not itself ensure this alignment; the
perturbation bound quantifies how flag error can erode the ideal rate.

\begin{assumption}[Invariant Hessian blocks]\label{ass:rigid-quadratic}
\begin{gather}
 F(w)=\tfrac12(w-w^\star)^\top H(w-w^\star),\quad H\succ0,\\
 H\Pi_j=\Pi_jH,\quad
 0<\ell_j\Pi_j\preceq\Pi_jH\Pi_j\preceq u_j\Pi_j,\quad
 B=\sum_{j=1}^{K+1}a_j\Pi_j,\quad a_j>0.
\end{gather}
The projectors in this assumption act on the parameter vector; a matrix
right-action uses $I_m\otimes\Pi_j$ after vectorization.
\end{assumption}
\begin{theorem}[Block conditioning and optimal scalar weights]\label{thm:rigid-conditioning}
Under Assumption~\ref{ass:rigid-quadratic}, let
\begin{gather}
 \ell_B:=\min_j a_j\ell_j,\quad u_B:=\max_j a_ju_j,\quad
 \kappa_B:=u_B/\ell_B,\quad \eta_\star:=2/(u_B+\ell_B),\\
 w_{t+1}=w_t-\eta B\nabla F(w_t).
\end{gather}
Then
\begin{align}
 \|w_t-w^\star\|_{B^{-1}}
 &\le\left[\max_{\lambda\in[\ell_B,u_B]}|1-\eta\lambda|\right]^t
                  \|w_0-w^\star\|_{B^{-1}},\\
 F(w_t)&\le\left(\frac{\kappa_B-1}{\kappa_B+1}\right)^{2t}F(w_0)
                 \quad(\eta=\eta_\star),\\
 \inf_{a_1,\ldots,a_{K+1}>0}\frac{\max_j a_ju_j}{\min_j a_j\ell_j}
 &=\max_j\frac{u_j}{\ell_j},\qquad a_j=1/\ell_j\text{ attains the infimum}.
\end{align}
If $H=\sum_jh_j\Pi_j$ and $a_j=h_j^{-1}$, then $w_1=w^\star$ for $\eta=1$.
\end{theorem}
\begin{proof}
\begin{align}
 x_t&:=B^{-1/2}(w_t-w^\star),\quad
 \widetilde H:=B^{1/2}HB^{1/2}=BH,\\
 x_{t+1}&=(I-\eta\widetilde H)x_t,\quad
 \operatorname{spec}(\widetilde H)\subset[\ell_B,u_B],\\
 F(w_t)&=\tfrac12x_t^\top\widetilde Hx_t,\quad
 \widetilde H^{1/2}x_t=(I-\eta\widetilde H)^t\widetilde H^{1/2}x_0,\\
 \min_{\eta>0}\max\{|1-\eta\ell_B|,|1-\eta u_B|\}
 &=\frac{u_B-\ell_B}{u_B+\ell_B},\\
 \frac{\max_i a_iu_i}{\min_i a_i\ell_i}
 &\ge\frac{a_ju_j}{a_j\ell_j}=u_j/\ell_j\quad(\forall j),\\
 a_j=1/\ell_j&\Longrightarrow\min_j a_j\ell_j=1,\quad
                                  \max_j a_ju_j=\max_j u_j/\ell_j.
\end{align}
\end{proof}
\begin{proposition}[Flag perturbation of quadratic conditioning]\label{prop:rigid-quadratic-error}
Let $B_\star=B(U^\star)$, $B=B(U)$, and
$\delta_B:=\Gamma^a d_P(U,U^\star)$. If
$\ell_\star I\preceq H^{1/2}B_\star H^{1/2}\preceq u_\star I$, then
\begin{gather}
 \ell_\star-\|H\|_{\op}\delta_B
 \le\lambda_{\min}(H^{1/2}BH^{1/2})
 \le\lambda_{\max}(H^{1/2}BH^{1/2})
 \le u_\star+\|H\|_{\op}\delta_B,\\
 \|H\|_{\op}\delta_B<\ell_\star
 \Longrightarrow\quad
 \kappa(H^{1/2}BH^{1/2})
 \le\frac{u_\star+\|H\|_{\op}\delta_B}
            {\ell_\star-\|H\|_{\op}\delta_B}.
\end{gather}
\end{proposition}
\begin{proof}
\begin{equation}
 \|H^{1/2}(B-B_\star)H^{1/2}\|_{\op}
 \le\|H\|_{\op}\|B-B_\star\|_{\op}\le\|H\|_{\op}\delta_B.
\end{equation}
\end{proof}
The small-gain criterion addresses a different question: when does direct
parameter contraction dominate feedback through optimizer state? The resolvent
$(I-D)^{-1}$ sums repeated propagation inside the state subsystem. Its coupling
cost must remain below the direct contraction margin $\alpha$, giving an exact
criterion for the stated nonnegative envelope.

\begin{theorem}[Schur-resolvent small gain]\label{thm:rigid-schur}
This is the nonsingular $M$-matrix/Schur-complement criterion for a nonnegative
feedback system~\citep{berman1994nonnegative}, specialized to the optimizer's
parameter--state partition.
Let a nonnegative sensitivity envelope at a fixed step size have the form
\begin{gather}
 \mathcal A_\eta:=
 \begin{pmatrix}1-\eta\alpha&\eta b^\top\\c&D\end{pmatrix}\ge0,\quad
 b,c\ge0,\quad D\ge0,\quad\rho(D)<1,\quad\eta>0.
\end{gather}
Then
\begin{equation}
 \rho(\mathcal A_\eta)<1\quad\Longleftrightarrow\quad
                      \alpha>b^\top(I-D)^{-1}c.
\end{equation}
If the right side holds, choose $\varepsilon>0$ so small that
\begin{gather}
 v^\top:=(b^\top+\varepsilon\mathbf1^\top)(I-D)^{-1}>0,\quad
 \delta:=\alpha-v^\top c>0,\quad a_\eta:=(1,\eta v^\top),\\
 0<\gamma\le\delta,\quad
 \gamma\eta\max_jv_j\le\varepsilon.
\end{gather}
Then
\begin{equation}
 a_\eta\mathcal A_\eta\preceq(1-\gamma\eta)a_\eta.
\end{equation}
\end{theorem}
\begin{proof}
\begin{align}
 (I-D)^{-1}&=\sum_{k=0}^\infty D^k\ge0,\\
 v^\top(I-D)&=b^\top+\varepsilon\mathbf1^\top,\\
 a_\eta\mathcal A_\eta
 &=(1-\eta(\alpha-v^\top c),\ \eta v^\top-\eta\varepsilon\mathbf1^\top)
 \preceq(1-\gamma\eta)a_\eta.
\end{align}
Thus $\rho(\mathcal A_\eta)\le1-\gamma\eta<1$. Conversely,
$\rho(\mathcal A_\eta)<1$ implies the existence of $x=(x_0,x_s)>0$ with
$(I-\mathcal A_\eta)x>0$, e.g.
$x=(I-\mathcal A_\eta)^{-1}\mathbf1$. Hence
\begin{gather}
 (I-D)x_s>cx_0,\quad x_s>(I-D)^{-1}cx_0,\quad
 \alpha x_0>b^\top x_s\ge b^\top(I-D)^{-1}cx_0.
\end{gather}
\end{proof}
\begin{corollary}[A uniform $1/n$ sensitivity floor]\label{cor:rigid-schur-stability}
Suppose a coupled optimizer satisfies
\begin{gather}
 \binom{u_{t+1}}{s_{t+1}}\preceq
 \mathcal A_\eta\binom{u_t}{s_t}
 +\frac{I_t}{b_t}\binom{\eta d}{h},\quad
 \PP(I_t=1)=b_t/n,\quad (u_0,s_0)=0,
\end{gather}
with the hypotheses of Theorem~\ref{thm:rigid-schur}. Then
\begin{align}
 \EE[u_T+\eta v^\top s_T]
 &\le\frac{d+v^\top h}{n\gamma}
                         (1-(1-\gamma\eta)^T),\\
 \epsilon_T^{\mathrm{stab}}&\le\frac{L_0(d+v^\top h)}{n\gamma}.
\end{align}
\end{corollary}
\begin{proof}
\begin{gather}
 \EE[u_{t+1}+\eta v^\top s_{t+1}]
 \le(1-\gamma\eta)\EE[u_t+\eta v^\top s_t]
                                      +\eta(d+v^\top h)/n.
\end{gather}
\end{proof}
\begin{proposition}[Local Jacobian envelopes for the actual map]\label{prop:rigid-jacobian}
Let $\mathcal T_{t,S}$ denote the complete map of
Definition~\ref{def:rigid-map}, expressed in regular local coordinates. For a convex
coupling region $\mathcal K$, partition its Jacobian into parameter and state blocks:
\begin{gather}
 \mathrm D\mathcal T_{t,S}(x)=
 \begin{pmatrix}J_{ww}&J_{ws}\\J_{sw}&J_{ss}\end{pmatrix},\quad
 \sup_{x\in\mathcal K}\|J_{ww}\|\le1-\eta\alpha,\quad
 \sup_{x\in\mathcal K}\|J_{ws}\|\le\eta b,\\
 \sup_{x\in\mathcal K}\|J_{sw}\|\le c,\quad
 \sup_{x\in\mathcal K}\|J_{ss}\|\le q<1.
\end{gather}
If $\alpha(1-q)>bc$, the scalar-state version of
Theorem~\ref{thm:rigid-schur} applies. Block norms give the matrix-state version.
For the broad triangle-inequality envelope of
Definition~\ref{def:rigid-transition}, these contraction conditions are additional
hypotheses, not consequences of smoothness alone.
\end{proposition}
\begin{proof}
\begin{equation}
 \mathcal T_{t,S}(x)-\mathcal T_{t,S}(x')
 =\int_0^1\mathrm D\mathcal T_{t,S}(x'+a(x-x'))(x-x')\,\mathrm da.
\end{equation}
\end{proof}

\subsection{A second generalization route: a full-rank flag posterior}
Stability follows perturbations along the entire training trajectory. Here we
instead control a distribution around the final parameters, allowing a
generalization statement when products of state sensitivities are too large
to be informative. The posterior may use the learned flag because the bound
holds simultaneously over posteriors; the prior must remain independent of
the sample to which the bound is applied.

\begin{definition}[Prior, posterior, and binary relative entropy]
\begin{gather}
 0\le\ell(w,z)\le1,\quad
 L(Q):=\EE_{w\sim Q,z\sim\mathcal D}\ell(w,z),\quad
 \widehat L_S(Q):=\EE_{w\sim Q}\frac1n\sum_i\ell(w,z_i),\\
 \operatorname{kl}(a\|b):=a\log(a/b)+(1-a)\log((1-a)/(1-b)),\\
 \operatorname{KL}(Q\|P):=\int\log(\mathrm dQ/\mathrm dP)\,\mathrm dQ,
 \quad Q\ll P,\\
 \operatorname{kl}_+^{-1}(a,c)
 :=\sup\{b\in[a,1]:\operatorname{kl}(a\|b)\le c\},\quad c\ge0.
\end{gather}
The prior $P$ is independent of $S$; the posterior may depend on the complete
training trajectory.
\end{definition}
\begin{lemma}[Seeger--Maurer exponential moment]\label{lem:rigid-pac-moment}
For iid $X_i\in[0,1]$, $p:=\EE X_i$, and $\overline X=n^{-1}\sum_iX_i$,
\begin{gather}
 \EE\exp\{n\operatorname{kl}(\overline X\|p)\}\le\Xi_n,\qquad
 \Xi_n:=\sum_{k=0}^n\binom nk(k/n)^k(1-k/n)^{n-k},\\
 \Xi_n\le n+1\quad(n\ge1),\qquad
 \Xi_n\le2\sqrt n\quad(n\ge8).
\end{gather}
Here $0^0=1$. The moment bound is attained by every nondegenerate Bernoulli
law, so $\Xi_n$ is its exact uniform prefactor~\citep{maurer2004note}.
\end{lemma}
\begin{proof}
For $p\in(0,1)$, the function $f(x):=\exp(n\operatorname{kl}(x\|p))$ is convex.
For every convex $g$ and $X\in[0,1]$,
\begin{gather}
 g(X)\le(1-X)g(0)+Xg(1),\qquad
 \EE g(X)\le(1-p)g(0)+pg(1).
\end{gather}
Applying this replacement separately to each coordinate yields independent
$B_i\sim\operatorname{Bernoulli}(p)$ and
\begin{align}
 \EE f(\overline X)
 &\le\EE f(n^{-1}\textstyle\sum_iB_i)\\
 &=\sum_{k=0}^n\binom nk p^k(1-p)^{n-k}
             \left(\frac{k/n}{p}\right)^k
             \left(\frac{1-k/n}{1-p}\right)^{n-k}\\
 &=\sum_{k=0}^n\binom nk(k/n)^k(1-k/n)^{n-k}=\Xi_n\le n+1.
\end{align}
For Bernoulli $X_i$ the comparison is an equality. Maurer's Stirling bound
\citep[Theorem 1]{maurer2004note} gives $\Xi_n\le2\sqrt n$ for $n\ge8$.
At $p\in\{0,1\}$, $X_i=p$ almost surely and the left side equals $1$.
\end{proof}
\begin{theorem}[Seeger--Maurer PAC--Bayes bound, without Pinsker relaxation]
\label{thm:rigid-pac}
For any sample-independent $P$ and $\delta\in(0,1)$, with probability at least
$1-\delta$ over $S$, simultaneously for all $Q\ll P$, with
$k_Q:=\operatorname{KL}(Q\|P)+\log(\Xi_n/\delta)$,
\begin{gather}
 \operatorname{kl}(\widehat L_S(Q)\|L(Q))
 \le k_Q/n,\\
 L(Q)\le\operatorname{kl}_+^{-1}(\widehat L_S(Q),k_Q/n)
 \le\min\left\{1,\widehat L_S(Q)+\sqrt{\frac{k_Q}{2n}}\right\}.
\end{gather}
This is the binary-KL form of the PAC--Bayes theorem
\citep{seeger2002pac,maurer2004note}; the square-root expression is a secondary
relaxation rather than the reported certificate.
\end{theorem}
\begin{proof}
Fubini, Lemma~\ref{lem:rigid-pac-moment}, and Markov's inequality give
\begin{gather}
 \EE_S\EE_{w\sim P}e^{n\operatorname{kl}(\widehat L_S(w)\|L(w))}\le\Xi_n,\\
 \PP_S\left(\EE_{w\sim P}e^{n\operatorname{kl}(\widehat L_S(w)\|L(w))}
                       >\frac{\Xi_n}{\delta}\right)\le\delta.
\end{gather}
On the complementary event, for $f(w):=n\operatorname{kl}(\widehat L_S(w)\|L(w))$,
the Gibbs/Donsker--Varadhan change-of-measure inequality
\citep{donsker1975asymptotic} and joint convexity of binary relative entropy give
\begin{align}
 \EE_Q f-\operatorname{KL}(Q\|P)
 &\le\log\EE_Pe^f,\\
 n\operatorname{kl}(\widehat L_S(Q)\|L(Q))
 &\le\EE_Qn\operatorname{kl}(\widehat L_S(w)\|L(w))
 \le\operatorname{KL}(Q\|P)+\log(\Xi_n/\delta),\\
 \operatorname{kl}(a\|b)&\ge2(a-b)^2.
\end{align}
The last inequality follows from
$\partial_a^2\operatorname{kl}(a\|b)=1/[a(1-a)]\ge4$ and Taylor expansion at
$a=b$. All assertions extend by continuity at boundary values.
\end{proof}
\begin{corollary}[Exact Catoni-type exponential inversion]\label{cor:catoni-inverse}
For $a\in[0,1]$ and $c\ge0$,
\begin{equation}
 \operatorname{kl}_+^{-1}(a,c)
 =\inf_{\gamma>0}\frac{1-\exp(-\gamma a-c)}{1-\exp(-\gamma)}.
\end{equation}
Thus the preceding theorem can be evaluated by a one-dimensional monotone
root or by exponential-temperature optimization~\citep{catoni2007pac}.
In particular, $\widehat L_S(Q)=0$ gives the exact KL-event consequence
\begin{equation}
 L(Q)\le1-e^{-k_Q/n}\le k_Q/n,
\end{equation}
which retains the fast $k_Q/n$ dependence lost by a square-root relaxation.
\end{corollary}
\begin{proof}
For $b\ge a$, Bernoulli entropy duality gives
\begin{equation}
 \operatorname{kl}(a\|b)
 =\sup_{\gamma\ge0}\{-\gamma a-\log(1-b+be^{-\gamma})\}.
\end{equation}
The constraint $\operatorname{kl}(a\|b)\le c$ is therefore equivalent to
$b\le(1-e^{-\gamma a-c})/(1-e^{-\gamma})$ for every $\gamma>0$.
Taking the largest admissible $b$ proves the identity, including boundary
values by limits. This optimization is algebraic on one uniform KL event and
incurs no additional selection penalty.
\end{proof}
\begin{proposition}[Exact flag-Gaussian complexity]\label{prop:rigid-pac-flag}
For $W\in\RR^{m\times d}$, let
\begin{gather}
 P=\mathcal N(0,\sigma_0^2I_{md}),\quad
 Q_{W,U,\sigma}=\mathcal N(\operatorname{vec}W,I_m\otimes\Sigma_U),\\
 \Sigma_U:=\sum_{j=1}^{K+1}\sigma_j^2\Pi_j(U),\quad
 \sigma_j>0,\quad\sigma_0>0.
\end{gather}
Using rowwise vectorization,
\begin{align}
 \operatorname{KL}(Q_{W,U,\sigma}\|P)
 &=\frac{\|W\|_F^2}{2\sigma_0^2}
 +\frac m2\sum_{j=1}^{K+1}q_j
       \left[\frac{\sigma_j^2}{\sigma_0^2}-1
                         -\log\frac{\sigma_j^2}{\sigma_0^2}\right],\\
 \tr(I_m\otimes\Sigma_U)&=m\sum_{j=1}^{K+1}q_j\sigma_j^2,\qquad
 \det(I_m\otimes\Sigma_U)=\prod_{j=1}^{K+1}(\sigma_j^2)^{mq_j}.
\end{align}
For $q_{K+1}=d-R>0$,
\begin{equation}
 \sigma_{K+1}\downarrow0\quad\Longrightarrow\quad
 \operatorname{KL}(Q_{W,U,\sigma}\|P)\longrightarrow\infty.
\end{equation}
\end{proposition}
\begin{proof}
For positive definite $\Sigma\in\RR^{D\times D}$ and $Q=\mathcal N(\mu,\Sigma)$,
\begin{align}
 \log\frac{\mathrm dQ}{\mathrm dP}(x)
 &=\tfrac12\log\frac{\det(\sigma_0^2I_D)}{\det\Sigma}
 -\tfrac12(x-\mu)^\top\Sigma^{-1}(x-\mu)
 +\frac{\|x\|^2}{2\sigma_0^2},\\
 \EE_Q(x-\mu)^\top\Sigma^{-1}(x-\mu)&=D,\qquad
 \EE_Q\|x\|^2=\|\mu\|^2+\tr\Sigma,\\
 \log\det\Sigma_U&=\sum_jq_j\log\sigma_j^2,\quad
 \tr\Sigma_U=\sum_jq_j\sigma_j^2.
\end{align}
\end{proof}
Positive residual variance keeps the posterior full rank and gives finite
complexity relative to the full-rank Gaussian prior. Under this isotropic
prior, the displayed KL depends on block dimensions and variances, while
posterior orientation can affect the perturbed loss. To obtain a statement
about the trained parameter $W$, the next result controls both empirical and
population changes induced by the Gaussian perturbation.

\begin{corollary}[A deterministic-center bound]\label{cor:rigid-pac-center}
In addition to $0\le\ell\le1$, suppose $\ell(\cdot,z)$ has
$L_\ell$-Lipschitz gradient for every $z$. On the event of
Theorem~\ref{thm:rigid-pac}, simultaneously for all data-dependent
$W,U,\sigma_j>0$,
set $v_\sigma:=m\sum_jq_j\sigma_j^2$,
$a_Q:=\min\{1,F_S(W)+L_\ell v_\sigma/2\}$. Then
\begin{equation}
 F_\mathcal D(W)\le\min\left\{1,\frac{L_\ell v_\sigma}{2}
           +\operatorname{kl}_+^{-1}(a_Q,k_Q/n)\right\},
\end{equation}
where $k_Q$ uses the exact Gaussian KL above. In particular,
\begin{align}
 F_\mathcal D(W)
 &\le F_S(W)+L_\ell m\sum_jq_j\sigma_j^2\\
 &\quad+\sqrt{\frac{\|W\|_F^2/(2\sigma_0^2)
 +\frac m2\sum_jq_j[\sigma_j^2/\sigma_0^2-1-\log(\sigma_j^2/\sigma_0^2)]
 +\log(\Xi_n/\delta)}{2n}}.
\end{align}
\end{corollary}
\begin{proof}
For $E\sim\mathcal N(0,I_m\otimes\Sigma_U)$,
\begin{gather}
 |\ell(W+E,z)-\ell(W,z)-\langle\nabla\ell(W,z),E\rangle|
                  \le\tfrac12L_\ell\|E\|_F^2,\\
 \EE E=0,\quad\EE\|E\|_F^2=m\sum_jq_j\sigma_j^2,\\
 |L(Q)-F_\mathcal D(W)|\le\tfrac12L_\ell m\sum_jq_j\sigma_j^2,\quad
 |\widehat L_S(Q)-F_S(W)|\le\tfrac12L_\ell m\sum_jq_j\sigma_j^2.
\end{gather}
Insert these two Taylor bounds into the inverse-KL certificate; monotonicity
of $\operatorname{kl}_+^{-1}$ in its first argument gives the first assertion.
Applying Pinsker only afterward gives the displayed analytic relaxation.
\end{proof}

\subsection{Convergence of the actual current-gradient-normalized update}
The second-moment update includes the current stochastic gradient, so its
normalizer and gradient are dependent. We analyze their combined nonlinear
field. The proof first establishes monotonicity of that field, then compares
the conditional noise distribution with its reflection. This provides a
descent term together with a quantitative cost for gradient bias and noise
asymmetry, before momentum and refresh errors are introduced.

\begin{definition}[Conditional normalization field]\label{def:actual-field}
Condition on $\mathcal F_{t-1}$, with $U_t$, the stored moments, and their ages
fixed. Use rowwise vectorization, $N=md$, and let $\mathcal P_{k,t}$ be the
orthogonal parameter-coordinate projector induced by $\Pi_{k,t}$.
For each active coordinate $i$ of block $k$ set
\begin{gather}
 \alpha_{k,t}:=\frac{1-\beta_{2k}}{1-\beta_{2k}^{c_{V,k,t}}},\quad
 c_{i,t}:=\frac{\beta_{2k}V_{i,t-1}^+}{1-\beta_{2k}^{c_{V,k,t}}},\quad
 h_{i,t}(z):=\frac{a_k}{\sqrt{c_{i,t}+\alpha_{k,t}z^2}+\epsilon},\\
 \mathcal B_t(g):=\bigoplus_{k=1}^K\diag_i h_{i,t}((\operatorname{mat}(g)U_{k,t})_i)
 \ \oplus\ \frac{a_\perp}{\|\operatorname{mat}(g)\Pi_{K+1,t}\|/\sqrt N+\epsilon}I_{m(d-R)},\\
 \mathcal T_t(g):=\mathcal B_t(g)g.
\end{gather}
The direct sum is expressed in the orthonormal block frame and lifted back to
parameter space. In particular, $\mathcal B_t(g)$ depends on the current $g$.
\end{definition}
\begin{lemma}[Strong monotonicity of the normalization field]\label{lem:actual-normalization}
Suppose $\|g_s\|\le G$ for all preceding steps, $a_k,a_\perp>0$, and
\begin{gather}
 M:=\epsilon^{-1}\max\{a_1,\ldots,a_K,a_\perp\},\\
 \kappa:=\min\left\{\min_k\frac{a_k\epsilon}{(3G+\epsilon)^2},
                 \frac{a_\perp\epsilon}{(3G/\sqrt N+\epsilon)^2}\right\}>0.
\end{gather}
Then, on $\|x\|,\|y\|\le3G$,
\begin{gather}
 \mathcal T_t(-x)=-\mathcal T_t(x),\quad
 \|\mathcal T_t(x)-\mathcal T_t(y)\|\le M\|x-y\|,\\
 \langle x-y,\mathcal T_t(x)-\mathcal T_t(y)\rangle
                              \ge\kappa\|x-y\|^2.
\end{gather}
\end{lemma}
\begin{proof}
The exact reset-age recursion gives
\begin{align}
 V_{i,t-1}^+&\le(1-\beta_{2k}^{c_{V,k,t}-1})G^2,\\
 0\le c_{i,t}&\le(1-\alpha_{k,t})G^2,\quad0<\alpha_{k,t}\le1,\\
 s(z)&:=\sqrt{c_{i,t}+\alpha_{k,t}z^2},\quad
 f(z):=a_kz/(s(z)+\epsilon),\\
 f'(z)&=\frac{a_k(c_{i,t}/s(z)+\epsilon)}{(s(z)+\epsilon)^2}
 \quad(s(z)>0),\\
 |z|\le3G&\Longrightarrow s(z)\le3G,\quad
 \frac{a_k\epsilon}{(3G+\epsilon)^2}\le f'(z)\le a_k/\epsilon.
\end{align}
For $c_{i,t}=z=0$, $f'(0)=a_k/\epsilon$ by continuity. On the residual block,
Proposition~\ref{prop:rigid-radial} gives the same bounds with
$3G$ replaced by $3G/\sqrt N$. Thus
\begin{gather}
 \kappa I\preceq\mathrm D\mathcal T_t(x)\preceq MI,\quad\|x\|\le3G,\\
 \mathcal T_t(x)-\mathcal T_t(y)
 =\int_0^1\mathrm D\mathcal T_t(y+s(x-y))(x-y)\,\mathrm ds.
\end{gather}
\end{proof}
\begin{definition}[Bias and symmetry defect]\label{def:actual-symmetry}
\begin{gather}
 h_t:=\nabla F(w_t),\quad\overline h_t:=\EE_tg_t,\quad
 b_t:=\overline h_t-h_t,\quad\xi_t:=g_t-\overline h_t,\\
 \nu_t:=\mathcal L(\xi_t\mid\mathcal F_{t-1}),\quad
 \nu_t^-:=(-\Id)_\#\nu_t,\quad
 \delta_t:=W_1(\nu_t,\nu_t^-),\quad D_t:=\|b_t\|+\delta_t/2,\\
 \mathfrak a_t:=\|\EE_t\mathcal T_t(\xi_t)\|,\qquad
 \widehat D_t:=\|b_t\|+\mathfrak a_t/M,\\
 W_1(\mu,\nu):=\inf_{(X,Y):\,X\sim\mu,Y\sim\nu}\EE\|X-Y\|.
\end{gather}
\end{definition}
\begin{theorem}[Odd-field coercivity and Kantorovich--Rubinstein calibration]
\label{thm:actual-coercivity}
If $\|g_t\|,\|h_t\|\le G$, then
\begin{align}
 \langle h_t,\EE_t[\mathcal B_t(g_t)g_t]\rangle
 &\ge\kappa\|h_t\|^2-M\widehat D_t\|h_t\|\\
 &\ge\tfrac\kappa2\|h_t\|^2-\frac{M^2}{2\kappa}\widehat D_t^2.
\end{align}
Moreover,
\begin{gather}
 \mathfrak a_t\le\min\{M\delta_t/2,\mathsf H_t\},\quad
 \widehat D_t\le D_t,\\
 \mathsf H_t^2:=m\sum_{k=1}^K\frac{q_ka_k^2}{\alpha_{k,t}}+Na_\perp^2.
\end{gather}
If $b_t=0$ and $\nu_t=\nu_t^-$, the stronger lower bound is
$\kappa\|h_t\|^2$; more generally $b_t=0$ and $\mathfrak a_t=0$ suffice.
\end{theorem}
\begin{proof}
Strong monotonicity applies on the segment from $\xi_t$ to $\xi_t+h_t$,
whose norm is at most $3G$. Lipschitz continuity controls the shift $b_t$:
\begin{align}
 \|\EE_t[\mathcal T_t(g_t)-\mathcal T_t(h_t+\xi_t)]\|
 &\le M\|b_t\|,\\
 \langle h_t,\EE_t[\mathcal T_t(h_t+\xi_t)-\mathcal T_t(\xi_t)]\rangle
 &\ge\kappa\|h_t\|^2,\\
 M\widehat D_t\|h_t\|&\le\tfrac\kappa2\|h_t\|^2
                              +\tfrac{M^2}{2\kappa}\widehat D_t^2.
\end{align}
For a unit vector $v$, $v^\top\mathcal T_t$ is $M$-Lipschitz and odd.
Kantorovich--Rubinstein duality~\citep{villani2009optimal} therefore gives
\begin{equation}
 2\mathfrak a_t
 =\left\|\int\mathcal T_t\,\mathrm d\nu_t-
                  \int\mathcal T_t\,\mathrm d\nu_t^-\right\|
 \le M W_1(\nu_t,\nu_t^-).
\end{equation}
Each active normalized coordinate has magnitude at most
$a_k/\sqrt{\alpha_{k,t}}$, and the residual norm is at most
$a_\perp\sqrt N$. Orthogonality gives $\|\mathcal T_t(\xi_t)\|\le\mathsf H_t$,
which supplies the second calibration without a factor $\epsilon^{-1}$.
\end{proof}
Bias and asymmetry enter through different mechanisms. Even an unbiased
gradient can acquire a nonzero mean update under nonlinear normalization if
its noise is asymmetric. Conversely, conditional symmetry cancels the odd
normalization field at zero signal even when the noise has substantial
variance. The directly normalized defect $\mathfrak a_t$ can be smaller than
the Wasserstein envelope $M\delta_t/2$; the latter is a convenient sufficient
bound, while the saturation calibration avoids introducing an unnecessary
$\epsilon^{-1}$ factor into the asymmetry numerator.

\begin{proposition}[Exact ambient momentum and refresh defect]\label{prop:actual-momentum}
Write $G_t:=\operatorname{mat}(g_t)$. Let $\widehat M_{k,t}$ denote the bias-corrected first moment before refresh,
$A_{k,t}:=\widehat M_{k,t}U_{k,t}^\top$, and
\begin{gather}
 m_t:=\operatorname{vec}\left(\sum_k A_{k,t}+G_t\Pi_{K+1,t}\right),\quad
 \theta_{k,t}:=\frac{\beta_{1k}(1-\beta_{1k}^{t-1})}{1-\beta_{1k}^t},\quad
 \Theta_t:=\sum_k\theta_{k,t}\mathcal P_{k,t},\\
 \beta:=\max_k\beta_{1k}<1,\quad G_m:=\sqrt{K+1}G.
\end{gather}
With the first-moment ages $c_{M,k,t}=t$, exactly,
\begin{gather}
 m_t=\Theta_tm_{t-1}+(I-\Theta_t)g_t+u_t,\qquad
 d_t=\mathcal B_t(g_t)m_t,\qquad \|m_t\|\le G_m,\\
 u_t=\operatorname{vec}\sum_k\theta_{k,t}
 \left[\widehat M_{k,t-1}O_{k,t-1}^\top U_{k,t}^\top
                          -\operatorname{mat}(m_{t-1})\Pi_{k,t}\right].
\end{gather}
Here $O_{k,t-1}=I$ without a preceding refresh. Furthermore,
\begin{align}
 \|u_t\|
 &\le\beta\sum_k\left[
 G\|O_{k,t-1}^\top-U_{k,t-1}^\top U_{k,t}\|_{\op}
 +G_m\|\Pi_{k,t}-\Pi_{k,t-1}\|_{\op}\right],\\
 \chi_{t-1}=0&\Longrightarrow u_t=0.
\end{align}
\end{proposition}
\begin{proof}
\begin{align}
 \widehat M_{k,t}
 &=\theta_{k,t}\widehat M_{k,t-1}O_{k,t-1}^\top
                         +(1-\theta_{k,t})G_tU_{k,t},\\
 \|m_t\|^2&=\sum_k\|\widehat M_{k,t}\|_F^2+\|G_t\Pi_{K+1,t}\|_F^2
 \le(K+1)G^2,\\
 A_{k,t-1}(I-\Pi_{k,t-1})&=0,\quad
 (\operatorname{mat}(m_{t-1})-A_{k,t-1})\Pi_{k,t-1}=0,\\
 \widehat M_{k,t-1}O^\top U_{k,t}^\top-A_{k,t-1}\Pi_{k,t}
 &=\widehat M_{k,t-1}(O^\top-U_{k,t-1}^\top U_{k,t})U_{k,t}^\top,\\
 (\operatorname{mat}(m_{t-1})-A_{k,t-1})\Pi_{k,t}
 &=\operatorname{mat}(m_{t-1})(I-\Pi_{k,t-1})
                                     (\Pi_{k,t}-\Pi_{k,t-1}).
\end{align}
\end{proof}
\begin{theorem}[Actual Adam normalization, momentum, and flag transport]
\label{thm:actual-convergence}
Assume $F$ is $L$-smooth, lower bounded, $\|h_t\|,\|g_t\|\le G$, $\lambda=0$,
and the exact map of Definition~\ref{def:rigid-map} is used. Define
\begin{gather}
 B_t:=\mathcal B_t(g_t),\quad
 C_t:=\eta_tB_t(I-\Theta_t)^{-1}\Theta_t,\quad C_0:=0\text{ before step }1,\\
 z_t:=w_t-C_{t-1}m_{t-1},\quad
 r_t:=(C_{t-1}-C_t)m_{t-1}-\eta_tB_t(I-\Theta_t)^{-1}u_t,\\
 \overline\eta:=\sup_t\eta_t,\quad
 G_z:=G+L\overline\eta M\beta G_m/(1-\beta),\quad
 S_T:=\sum_{t=1}^T\eta_t,\quad\mathsf J_t:=\min\{MG,\mathsf H_t\}.
\end{gather}
Then
\begin{align}
 \frac\kappa2\sum_{t=1}^T\eta_t\EE\|\nabla F(w_t)\|^2
 &\le\Delta_0+L\sum_{t=1}^T\eta_t^2\EE\mathsf J_t^2
       +LG_m\sum_{t=1}^T\eta_t\EE[\mathsf J_t\|C_{t-1}\|_{\op}]\\
 &\quad+\frac{M^2}{2\kappa}\sum_{t=1}^T\eta_t\EE \widehat D_t^2
       +G_z\sum_{t=1}^T\EE\|r_t\|+L\sum_{t=1}^T\EE\|r_t\|^2,
 \label{eq:actual-main-rate}\\
 \|r_t\|&\le G_m\|C_t-C_{t-1}\|_{\op}
                      +\frac{\eta_tM}{1-\beta}\|u_t\|.
\end{align}
Here $\Delta_0:=\EE F(w_1)-F_{\inf}$. No predictability assumption on $B_t$ is used.
\end{theorem}
\begin{proof}
The virtual-iterate identity is algebraic, hence remains valid for random $B_t$:
\begin{align}
 z_{t+1}&=z_t-\eta_tB_tg_t+r_t,\\
 \|w_t-z_t\|&\le G_m\|C_{t-1}\|_{\op},\quad
 \|C_t\|_{\op}\le\eta_t M\beta/(1-\beta),\quad
 \|B_tg_t\|\le\mathsf J_t,\\
 F(z_{t+1})&\le F(z_t)-\eta_t\langle\nabla F(z_t),B_tg_t\rangle
                 +\langle\nabla F(z_t),r_t\rangle
                 +L\eta_t^2\|B_tg_t\|^2+L\|r_t\|^2,\\
 \left|\langle\nabla F(z_t)-h_t,B_tg_t\rangle\right|
 &\le LG_m\mathsf J_t\|C_{t-1}\|_{\op},\\
 \EE_t\langle h_t,B_tg_t\rangle
 &\ge\tfrac\kappa2\|h_t\|^2-\tfrac{M^2}{2\kappa}\widehat D_t^2,\\
 \|\nabla F(z_t)\|&\le\|h_t\|+L\|w_t-z_t\|\le G_z.
\end{align}
Take conditional expectations and sum, using $z_1=w_1$ and $F(z_{T+1})\ge F_{\inf}$.
\end{proof}
The bound now separates current-normalization bias, momentum lag, and the
variation/transport remainder $r_t$. The virtual-iterate identity remains exact
even though $B_t$ is random and depends on $g_t$; the descent argument uses
coercivity of the combined field. A rate follows only when the accumulated
state changes and noise defect satisfy the budgets specified next.

\begin{corollary}[Observable variation and batch-size conditions]\label{cor:actual-rate}
Set $\eta_t=\eta_0T^{-1/2}$ and define
\begin{gather}
 V_T^B:=\sum_{t=1}^T\|B_t-B_{t-1}\|_{\op},\quad
 V_T^\Theta:=\sum_{t=1}^T\|\Theta_t-\Theta_{t-1}\|_{\op},\quad
 U_T:=\sum_{t=1}^T\|u_t\|,\quad B_0=\Theta_0=0.
\end{gather}
Then
\begin{align}
 \sum_t\|C_t-C_{t-1}\|_{\op}
 &\le\eta\left[\frac\beta{1-\beta}V_T^B+
                             \frac{M}{(1-\beta)^2}V_T^\Theta\right],\\
 \EE[V_T^B+V_T^\Theta+U_T]=O(\sqrt T)
 &\Longrightarrow\quad
 \frac1T\sum_{t=1}^T\EE\|\nabla F(w_t)\|^2
 =O(T^{-1/2})+\frac{M^2}{\kappa^2T}\sum_t\EE \widehat D_t^2.
\end{align}
For an average of $b$ independent unbiased sample gradients with conditional
variance at most $\sigma^2$, without clipping,
\begin{equation}
 \widehat D_t^2\le\EE_t\|g_t-h_t\|^2\le\sigma^2/b.
\end{equation}
Consequently the displayed variation condition and $b\gtrsim\sqrt T$ give
$T^{-1/2}$ stationarity. With norm clipping, replace the last error by
$2\|b_t^{\mathrm{clip}}\|^2+2\sigma^2/b$.
\end{corollary}
\begin{proof}
\begin{align}
 (I-\Theta)^{-1}\Theta&=(I-\Theta)^{-1}-I,\\
 \|(I-\Theta)^{-1}-(I-\Theta')^{-1}\|_{\op}
 &\le(1-\beta)^{-2}\|\Theta-\Theta'\|_{\op},\\
 W_1(\nu_t,\nu_t^-)&\le2\EE_t\|\xi_t\|
 \le2\sqrt{\EE_t\|\xi_t\|^2},\\
 \|r_t\|&\le O(\eta),\quad
 \sum_t\EE\|r_t\|=O(\eta\sqrt T)=O(1),\\
 \sum_t\EE\|r_t\|^2&\le(\sup_t\|r_t\|)\sum_t\EE\|r_t\|=O(\eta).
\end{align}
Clipping is nonexpansive, so for independent copies $X,X'$,
\begin{gather}
 \operatorname{Var}(\operatorname{clip}(X))
 =\tfrac12\EE\|\operatorname{clip}(X)-\operatorname{clip}(X')\|^2
 \le\tfrac12\EE\|X-X'\|^2=\operatorname{Var}(X),\\
 \widehat D_t^2\le D_t^2=(\|b_t\|+\delta_t/2)^2
 \le2\|b_t\|^2+2\operatorname{Var}_t(g_t).
\end{gather}
\end{proof}
\begin{proposition}[An exact asymmetry-induced stationary bias]\label{prop:actual-skew-counterexample}
Consider a fixed one-coordinate flag branch with $\beta_1=\beta_2=0$, $a=1$,
$F(w)=w^2/2$, and unbiased additive noise
\begin{gather}
 \PP(\xi=3\sigma)=1/10,\quad\PP(\xi=-\sigma/3)=9/10,\quad
 \EE\xi=0,\quad\EE\xi^2=\sigma^2,\\
 (\xi_t)_{t\ge1}\text{ iid},\quad w_{t+1}=w_t-\eta_t\frac{w_t+\xi_t}{|w_t+\xi_t|+\epsilon}.
\end{gather}
Its expected update field has a unique root
\begin{gather}
 \Psi_\epsilon(w):=\frac1{10}\frac{w+3\sigma}{|w+3\sigma|+\epsilon}
 +\frac9{10}\frac{w-\sigma/3}{|w-\sigma/3|+\epsilon},\\
 w_\epsilon=\tfrac12\left[
 \sqrt{(8\sigma/3+5\epsilon/4)^2+4\sigma^2}
                      -(8\sigma/3+5\epsilon/4)\right]\in(0,\sigma/3),\\
 \lim_{\epsilon\downarrow0}w_\epsilon=\sigma/3>0.
\end{gather}
For deterministic $\eta_t>0$, $\sup_t\eta_t<\infty$,
$\sum_t\eta_t=\infty$, $\sum_t\eta_t^2<\infty$, the iterates satisfy
\begin{equation}
 w_t\longrightarrow w_\epsilon\quad\text{a.s.},\qquad
 \|\nabla F(w_t)\|^2\longrightarrow w_\epsilon^2>0\quad\text{a.s.}
\end{equation}
\end{proposition}
\begin{proof}
\begin{align}
 \Psi_\epsilon'(w)&>0,\quad\Psi_\epsilon(0)<0,\quad
                                 \Psi_\epsilon(\sigma/3)>0,\\
 \Psi_\epsilon(w)=0,\quad0<w<\sigma/3
 &\Longleftrightarrow
 (w+3\sigma)(w-\sigma/3)+\tfrac54\epsilon w=0.
\end{align}
The scalar branch embeds in the exact flag algorithm as
\begin{gather}
 W=\begin{pmatrix}w&0\\0&0\end{pmatrix},\quad U=e_1\in\St(2,1),\quad
 G=\begin{pmatrix}w+\xi&0\\0&0\end{pmatrix},\\
 G\Pi_\perp=0,\quad C_k=\diag(c_k,0),\quad
 \operatorname{grad}\Phi(U;C)=0.
\end{gather}
If $w>\sigma/3$, every update decreases $w$; if $w<-3\sigma$, every update
increases $w$. Since every increment has magnitude at most $\eta_t$, all iterates
lie in a deterministic compact interval $\mathcal I$. On $\mathcal I$,
$\Psi_\epsilon'(w)\ge c_\mathcal I>0$, choosing
$c_\mathcal I\le(2\sup_t\eta_t)^{-1}$. Thus
\begin{align}
 \EE_t(w_{t+1}-w_\epsilon)^2
 &\le(w_t-w_\epsilon)^2-2\eta_t(w_t-w_\epsilon)\Psi_\epsilon(w_t)+\eta_t^2\\
 &\le(1-2c_\mathcal I\eta_t)(w_t-w_\epsilon)^2+\eta_t^2.
\end{align}
The almost-supermartingale theorem gives convergence of $(w_t-w_\epsilon)^2$
and $\sum_t\eta_t(w_t-w_\epsilon)^2<\infty$; hence its limit is zero.
\end{proof}
The counterexample converges to a zero of the expected normalized update field
that is separated from the minimizer of $F$. Its gradient noise is unbiased
and has bounded variance, so those properties alone cannot remove this
stationary bias. This gives a concrete reason to retain the symmetry-defect
term when interpreting the convergence bound for the actual algorithm.

\begin{proposition}[Resolving the variation budgets into actual state changes]
\label{prop:actual-budget-resolution}
For the exact preconditioner let $S_{k,t}:=\sqrt{\widehat V_{k,t}}$ and
$c_t:=(\|G_t\Pi_{K+1,t}\|_F/\sqrt N+\epsilon)^{-1}$. With constant $a_k$,
\begin{align}
 \|B_t-B_{t-1}\|_{\op}
 &\le\sum_k a_k\left[
 \epsilon^{-2}\|S_{k,t}-S_{k,t-1}\|_{\max}
 +2\epsilon^{-1}\|U_{k,t}-U_{k,t-1}\|_{\op}\right]\\
 &\quad+a_\perp|c_t-c_{t-1}|
       +(a_\perp/\epsilon)\|P_{K,t}-P_{K,t-1}\|_{\op},\\
 |c_t-c_{t-1}|
 &\le\frac{\|G_t-G_{t-1}\|_F+
                    G\|P_{K,t}-P_{K,t-1}\|_{\op}}{\epsilon^2\sqrt N},\\
 V_T^\Theta
 &\le\sum_k\beta_{1k}+\beta\sum_{t=2}^T\sum_k
                                  \|\Pi_{k,t}-\Pi_{k,t-1}\|_{\op}.
\end{align}
At refresh $s$ let $Z_s=\xi_s\operatorname{grad}\Phi(U_s;\widetilde C_s)$,
$z_s=\|Z_s\|_F$, and $e_s=(1+z_s^2)z_s^2<1$. The transport budget obeys
\begin{align}
 U_T&\le\beta\sqrt K\sum_{s\le S}\left[
 G\left(1+\frac2{2-e_s}\right)e_s
 +2G_m\left(z_s+e_s+\frac{2e_s}{2-e_s}\right)\right].
\end{align}
Consequently $z_s\le H\xi_0/\sqrt S$ and $S=O(T)$ imply
$U_T+V_T^\Theta=O(\sqrt T)$. This does not by itself control $V_T^B$:
second-moment resets and persistent gradient fluctuations remain explicit in
$\|S_{k,t}-S_{k,t-1}\|$ and $|c_t-c_{t-1}|$. A constant refresh step does not
imply the required variation rate.
\end{proposition}
\begin{proof}
Let $\mathcal U_{k,t}=I_m\otimes U_{k,t}$. Since
$B_{k,t}=a_k\mathcal U_{k,t}\diag((S_{k,t}+\epsilon)^{-1})\mathcal U_{k,t}^\top$,
\begin{gather}
 \|\diag((S+\epsilon)^{-1}-(S'+\epsilon)^{-1})\|_{\op}
 \le\epsilon^{-2}\|S-S'\|_{\max},\\
 \|\mathcal U\mathcal H\mathcal U^\top-
          \mathcal U'\mathcal H\mathcal U'^\top\|_{\op}
 \le2\|\mathcal H\|_{\op}\|U-U'\|_{\op}.
\end{gather}
Also $\theta_{k,t}\uparrow\beta_{1k}$, so its total scalar variation is
$\beta_{1k}$. Proposition~\ref{prop:rigid-small-transport} gives
\begin{gather}
 \|U^+-U\|_F\le z_s+e_s+2e_s/(2-e_s),\\
 \sum_k\|\Pi_k^+-\Pi_k\|_{\op}
 \le2\sqrt K\|U^+-U\|_F,\quad
 \sum_k\|O_k^\top-U_k^\top U_k^+\|_{\op}
 \le\sqrt K(1+2/(2-e_s))e_s.
\end{gather}
Insert these inequalities into Proposition~\ref{prop:actual-momentum}, and use
$\sum_sz_s=O(\sqrt S)$, $\sum_se_s=O(1)$.
\end{proof}
\begin{corollary}[Complete networks and dense fallback tensors]
\label{cor:actual-network}
For finitely many parameter blocks, apply the constructions blockwise and set
\begin{gather}
 B_t=\bigoplus_lB_{l,t},\quad\Theta_t=\bigoplus_l\Theta_{l,t},\quad
 m_t=(m_{1,t},\ldots,m_{L,t}),\quad u_t=(u_{1,t},\ldots,u_{L,t}),\\
 M=\max_lM_l,\quad\kappa=\min_l\kappa_l,\quad
 \beta=\max_l\beta_l,\quad G^2=\sum_lG_l^2,\quad G_m^2=\sum_lG_{m,l}^2.
\end{gather}
Set $\mathsf H_t^2:=\sum_l\mathsf H_{l,t}^2$ and
$\mathsf J_t:=\min\{MG,\mathsf H_t\}$ for the joint normalized field.
A dense fallback tensor is a fixed full-coordinate block with diagonal
$B_{l,t}$, scalar $\Theta_{l,t}$ and $u_{l,t}=0$. Theorem~\ref{thm:actual-convergence}
then holds for the joint, coupled network objective, using these aggregate
constants and a joint conditional symmetry defect. For decoupled weight decay
$w_{t+1}=(1-\eta_t\lambda)w_t-\eta_tB_tm_t$, its algebraic identity becomes
\begin{equation}
 z_{t+1}=z_t-\eta_tB_tg_t+\widetilde r_t,\qquad
 \widetilde r_t=r_t-\eta_t\lambda w_t.
\end{equation}
The finite-time inequality remains valid with $r_t$ replaced by
$\widetilde r_t$; decay is not silently absorbed into an unbiased data gradient.
\end{corollary}
\begin{proof}
\begin{gather}
 \|\bigoplus_lB_l\|_{\op}=\max_l\|B_l\|_{\op},\quad
 \sum_l\langle h_l,\mathcal T_l(h_l+\xi_l^s)-\mathcal T_l(\xi_l^s)\rangle
 \ge(\min_l\kappa_l)\sum_l\|h_l\|^2,\\
 \|m_t\|^2=\sum_l\|m_{l,t}\|^2\le\sum_lG_{m,l}^2.
\end{gather}
Apply the same virtual-iterate identity and smoothness inequality to the direct sum.
\end{proof}
\begin{lemma}[State-dependent curvature of the normalization field]
\label{lem:actual-mature-variance}
In Definition~\ref{def:actual-field}, write $a_i=a_{k(i)}$ and
$\alpha_{i,t}=\alpha_{k(i),t}$. Restrict active coordinate $i$ to
$|z_i|\le H_i$, and the residual to radius $H_\perp$. Set
\begin{gather}
 s_{i,\max}:=\sqrt{c_{i,t}+\alpha_{i,t}H_i^2},\qquad
 \kappa_{i,t}:=\frac{a_i(c_{i,t}/s_{i,\max}+\epsilon)}
                               {(s_{i,\max}+\epsilon)^2},\\
 \kappa_t:=\min\left\{\min_i\kappa_{i,t},
           \frac{a_\perp\epsilon}{(H_\perp/\sqrt N+\epsilon)^2}\right\}.
\end{gather}
The value at $s_{i,\max}=0$ is interpreted as $a_i/\epsilon$. These are exact
lower Jacobian bounds on the stated product domain. In particular,
\begin{gather}
 c_{i,t}\ge\nu^2>0,\quad s_{i,\max}\le S
 \quad\Longrightarrow\quad
 \kappa_{i,t}\ge\frac{a_i\nu^2}{S(S+\epsilon)^2},\\
 c_{i,t}=0\quad\Longrightarrow\quad
 \kappa_{i,t}=\frac{a_i\epsilon}{(\sqrt{\alpha_{i,t}}H_i+\epsilon)^2}.
\end{gather}
Thus a mature variance floor removes the explicit vanishing-$\epsilon$ factor
in an active block, whereas a reset restores it. The residual can still be the
limiting block. If the support domains cover $\xi_t+sh_t$, $0\le s\le1$,
Theorem~\ref{thm:actual-coercivity} holds with the predictable $\kappa_t$.
\end{lemma}
\begin{proof}
For $s>0$,
\begin{equation}
 \frac{\mathrm d}{\mathrm ds}\left(\frac{c/s+\epsilon}{(s+\epsilon)^2}\right)
 =-\frac{c}{s^2(s+\epsilon)^2}
   -\frac{2(c/s+\epsilon)}{(s+\epsilon)^3}<0.
\end{equation}
The active derivative is minimized at the maximal allowed $|z_i|$; the radial
residual derivative is minimized at its maximal radius. Apply the integral
Jacobian identity on the same convex product domain.
\end{proof}
The state-dependent lower bound explains why reset times deserve separate
attention. A positive historical variance can maintain curvature of an active
normalization field as $\epsilon$ decreases; immediately after resetting it,
the lower bound again depends explicitly on $\epsilon$. The residual may still
limit the joint constant, so this observation alone does not prescribe an
optimal regularization value.

\begin{corollary}[PL contraction for the actual normalized momentum map]
\label{cor:actual-PL}
Under Theorem~\ref{thm:actual-convergence}, additionally suppose, along the
actual iterates,
\begin{gather}
 \tfrac12\|\nabla F(w_t)\|^2\ge\mu(F(w_t)-F^\star),\quad
 0<\mu\kappa\eta_t\le1,\quad
 a_t:=\EE[F(z_t)-F^\star],\\
 A_t:=GG_m\|C_{t-1}\|_{\op}
                    +\tfrac L2G_m^2\|C_{t-1}\|_{\op}^2,\\
 Q_t:=\mu\kappa\eta_t\EE A_t+L\eta_t^2\EE\mathsf J_t^2
       +LG_m\eta_t\EE[\mathsf J_t\|C_{t-1}\|_{\op}]\\
 \hphantom{Q_t:=} +\frac{M^2\eta_t}{2\kappa}\EE \widehat D_t^2
                    +G_z\EE\|r_t\|+L\EE\|r_t\|^2.
\end{gather}
Then, with $q_t:=1-\mu\kappa\eta_t$,
\begin{align}
 a_{t+1}&\le q_ta_t+Q_t,\\
 a_{T+1}&\le\left(\prod_{t=1}^Tq_t\right)a_1
          +\sum_{t=1}^T\left(\prod_{s=t+1}^Tq_s\right)Q_t.
\end{align}
For constant $\eta$ and $Q_t\le K_0\eta^2+Z_0\eta$,
\begin{equation}
 a_{T+1}\le(1-\mu\kappa\eta)^Ta_1+
              \frac{K_0\eta+Z_0}{\mu\kappa}
                               [1-(1-\mu\kappa\eta)^T].
\end{equation}
The forcing explicitly retains bias/symmetry, moment normalization, and
transport/reset variation; the stated inequality on $Q_t$ is an additional
condition on these quantities.
\end{corollary}
\begin{proof}
\begin{align}
 |F(w_t)-F(z_t)|
 &\le\|\nabla F(w_t)\|\|w_t-z_t\|+\tfrac L2\|w_t-z_t\|^2\le A_t,\\
 -\tfrac\kappa2\eta_t\|\nabla F(w_t)\|^2
 &\le-\mu\kappa\eta_t(F(w_t)-F^\star)\\
 &\le-\mu\kappa\eta_t(F(z_t)-F^\star)+\mu\kappa\eta_t A_t.
\end{align}
Insert this in the one-step inequality in the proof of
Theorem~\ref{thm:actual-convergence} and expand the scalar recurrence.
\end{proof}

\subsection{Directional PAC--Bayes complexity and optimal flag variances}
The previous Gaussian bound uses one curvature constant for all directions.
An anisotropic envelope instead charges perturbations according to the
curvature of each flag block. We optimize the variance within every block,
then examine the effect of refining the flag and of tracking its oracle
inexactly. These are comparisons of posterior bounds under the stated
evaluation-loss assumptions.

\begin{assumption}[Anisotropic curvature envelope]\label{ass:directional-curvature}
For the bounded evaluation loss $0\le\ell(w,z)\le1$, assume a sample-independent
matrix $H_0\succeq0$ such that, globally in $w,z$,
\begin{equation}
 -H_0\preceq\nabla_w^2\ell(w,z)\preceq H_0.
\end{equation}
For $\mathcal P_j:=I_m\otimes\Pi_j(U)$, set
\begin{gather}
 d_j:=mq_j,\quad C_j(U):=\tr(\mathcal P_jH_0),\quad
 \overline h_j(U):=C_j(U)/d_j,\quad s_j:=\sigma_j^2>0,\\
 \Sigma_U:=\sum_j s_j\Pi_j(U),\quad
 \widetilde\Sigma_U:=I_m\otimes\Sigma_U=\sum_j s_j\mathcal P_j,\quad
 c_{n,\delta}:=\log(\Xi_n/\delta).
\end{gather}
The curvature assumption concerns the bounded evaluation loss; it does not
assert that unbounded training cross-entropy lies in $[0,1]$.
\end{assumption}
\begin{theorem}[Optimized directional posterior]\label{thm:directional-pac}
Under Assumption~\ref{ass:directional-curvature}, fix
$P=\mathcal N(0,\sigma_0^2I_N)$ independently of the sample. With probability
at least $1-\delta$, simultaneously for every learned $W,U$ and every $\lambda>0$,
\begin{align}
 F_\mathcal D(W)
 &\le F_S(W)+\frac{\|W\|_F^2}{2\lambda\sigma_0^2}
 +\frac{c_{n,\delta}}\lambda+\frac\lambda{8n}
 +\mathcal C_\lambda(U),\\
 \mathcal C_\lambda(U)
 &:=\frac1{2\lambda}\sum_{j=1}^{K+1}d_j
                 \log(1+2\lambda\sigma_0^2\overline h_j(U)).
 \label{eq:directional-complexity}
\end{align}
The optimal positive posterior variances in this bound are exactly
\begin{equation}
 s_j^\star(U,\lambda)
 =\frac{\sigma_0^2}{1+2\lambda\sigma_0^2\overline h_j(U)}.
\end{equation}
\end{theorem}
\begin{proof}
For $E\sim\mathcal N(0,\widetilde\Sigma_U)$, Taylor's formula gives
\begin{align}
 |\ell(W+E,z)-\ell(W,z)-\langle\nabla\ell(W,z),E\rangle|
 &\le\tfrac12 E^\top H_0E,\\
 |L(Q)-F_\mathcal D(W)|,\quad|\widehat L_S(Q)-F_S(W)|
 &\le\tfrac12\tr(H_0\widetilde\Sigma_U)=\tfrac12\sum_j C_js_j.
\end{align}
On the single event of Theorem~\ref{thm:rigid-pac},
\begin{align}
 F_\mathcal D(W)
 &\le F_S(W)+\sum_jC_js_j+
 \sqrt{\frac{\|W\|_F^2/(2\sigma_0^2)
   +\frac12\sum_jd_j[s_j/\sigma_0^2-1-\log(s_j/\sigma_0^2)]
   +c_{n,\delta}}{2n}},\\
 \sqrt{a/(2n)}&\le a/\lambda+\lambda/(8n)\quad(a\ge0,\lambda>0).
\end{align}
For each $j$, define the strictly convex scalar objective
\begin{gather}
 J_j(s):=C_js+\frac{d_j}{2\lambda}
                       [s/\sigma_0^2-1-\log(s/\sigma_0^2)],\\
 J_j'(s)=C_j+\frac{d_j}{2\lambda}(\sigma_0^{-2}-s^{-1}),\quad
 J_j''(s)=\frac{d_j}{2\lambda s^2}>0,\\
 J_j'(s_j^\star)=0,\quad
 J_j(s_j^\star)=\frac{d_j}{2\lambda}
                                  \log(1+2\lambda\sigma_0^2C_j/d_j).
\end{gather}
The event is uniform over all posteriors; optimization over data-dependent
$U,s_j,\lambda$ therefore requires no additional union bound.
\end{proof}
The optimal variance is smaller in blocks with greater average curvature:
perturbations there have a larger loss cost. The logarithmic complexity results
from balancing that cost against the KL penalty for shrinking variance.
Here $\lambda$ is a bound-optimization parameter; it is separate from the
optimizer's learning rate and from weight-decay notation used earlier.
The closed form optimizes this analytic envelope; an evaluated certificate
can be tighter by using the inverse-KL bound before its scalar relaxation.

\begin{proposition}[A strict benefit of posterior flag refinement]
\label{prop:directional-refinement}
Suppose a block $\mathcal P$ of dimension $d$ is refined into orthogonal blocks
$\mathcal P_1,\ldots,\mathcal P_J$, with
\begin{gather}
 \mathcal P=\sum_{j=1}^J\mathcal P_j,\quad d=\sum_jd_j,\quad
 h_j:=\tr(\mathcal P_jH_0)/d_j,\quad \overline h:=\sum_j(d_j/d)h_j.
\end{gather}
Then
\begin{equation}
 \sum_j\frac{d_j}{2\lambda}\log(1+2\lambda\sigma_0^2h_j)
 \le\frac d{2\lambda}\log(1+2\lambda\sigma_0^2\overline h).
\end{equation}
For $d_j>0$, equality holds if and only if $h_1=\cdots=h_J$. Thus refining a
fixed flag can strictly improve the optimized posterior bound when the block
curvatures differ. This concerns the posterior family, not an ordering of the
trained optimizers' test errors.
More quantitatively, let $h_-:=\min_jh_j$, $h_+:=\max_jh_j$,
$V_h:=\sum_j(d_j/d)(h_j-\overline h)^2$, and let $\mathcal G_\lambda$ be
the coarse cost minus the refined cost. Then
\begin{equation}
 \frac{d\lambda\sigma_0^4 V_h}{(1+2\lambda\sigma_0^2h_+)^2}
 \le\mathcal G_\lambda\le
 \frac{d\lambda\sigma_0^4 V_h}{(1+2\lambda\sigma_0^2h_-)^2}.
\end{equation}
The two sides agree to leading order as the block curvatures coalesce.
\end{proposition}
\begin{proof}
\begin{gather}
 f(h):=\log(1+2\lambda\sigma_0^2h),\qquad
 f''(h)=-\frac{(2\lambda\sigma_0^2)^2}{(1+2\lambda\sigma_0^2h)^2}<0,\\
 \sum_j(d_j/d)f(h_j)\le f(\sum_j(d_j/d)h_j).
\end{gather}
For the quantitative statement apply Taylor's formula at $\overline h$.
The weighted linear term cancels, and $-f''$ lies between its values at
$h_+$ and $h_-$. Multiplication by $d/(2\lambda)$ gives the two-sided
strong-concavity remainder. Continuity of $f''$ proves the local sharpness.
\end{proof}
Refinement permits different variances where a coarser posterior required a
shared one. Strict concavity makes this flexibility useful precisely when the
subblocks have unequal average curvature. Whether a covariance-trained flag
discovers such directions is a separate alignment question. The next result
controls the change in complexity caused by tracking a chosen oracle
inexactly, without assuming that oracle minimizes the curvature cost.

\begin{proposition}[Tracking error in the directional complexity]
\label{prop:directional-tracking}
Partition $H_0\in\RR^{md\times md}$ into $d\times d$ row-coordinate blocks
$H_{ab}$, and define $A_0:=\sum_{a=1}^m H_{aa}\succeq0$. Then
Set $s_A:=\lambda_{\max}(A_0)-\lambda_{\min}(A_0)$. Then
\begin{gather}
 C_j(U)=\tr(\Pi_j(U)A_0),\\
 |\mathcal C_\lambda(U)-\mathcal C_\lambda(V)|
 \le\sigma_0^2 s_A\sqrt d\,d_P(U,V).
\end{gather}
In the regular oracle chart, let $\widehat x_t^\star$ and $x_t^\star$ be the
oracles for $\widehat C_t$ and $C_t$, respectively, and $U_t=\psi(x_t)$. Then
\begin{align}
 |\mathcal C_\lambda(U_t)-\mathcal C_\lambda(\psi(x_t^\star))|
 &\le\sigma_0^2 s_A\sqrt d\,J_P
 \left(\|x_t-\widehat x_t^\star\|+
     \frac{2J_\psi}{\mu_\phi}\sum_k\omega_k\sqrt{r_k}\epsilon_{k,t}\right).
\end{align}
\end{proposition}
\begin{proof}
Centering by $c_A:=(\lambda_{\max}(A_0)+\lambda_{\min}(A_0))/2$ costs
nothing because every projector difference is traceless. The trace/nuclear
norm duality and the exact overlap identity give
For $g_j(c):=d_j(2\lambda)^{-1}\log(1+2\lambda\sigma_0^2c/d_j)$,
\begin{align}
 0\le g_j'(c)&=\frac{\sigma_0^2}{1+2\lambda\sigma_0^2c/d_j}\le\sigma_0^2,\\
 |C_j(U)-C_j(V)|
 &\le\|A_0-c_AI\|_{\op}\|\Pi_j(U)-\Pi_j(V)\|_*\\
 &\le\tfrac{s_A}{2}\sqrt{2q_j}\|\Pi_j(U)-\Pi_j(V)\|_F,\\
 \sum_{j=1}^{K+1}\|\Pi_j(U)-\Pi_j(V)\|_F^2
 &\le2d_P(U,V)^2,\\
 \sum_j\sqrt{2q_j}\|\Pi_j(U)-\Pi_j(V)\|_F
 &\le\sqrt{2\sum_jq_j}\sqrt{2d_P(U,V)^2}=2\sqrt d\,d_P(U,V).
\end{align}
Apply the triangle inequality and Lemma~\ref{lem:tf-oracle} to the two oracle branches.
In particular $A_0=cI$ gives zero orientation cost exactly. The bound depends
on spectral spread, removing an isotropic component that cannot distinguish flags.
\end{proof}
\begin{proposition}[A sample-independent anisotropic flag prior]
\label{prop:directional-prior}
Let a fixed prior flag $V$ have projectors $\Pi_l^0$ and ranks $q_l^0$, and let
\begin{gather}
 \Lambda_0:=\sum_l\tau_l^2\Pi_l^0\succ0,\quad
 P_V:=\mathcal N(0,I_m\otimes\Lambda_0),\quad
 Q_U:=\mathcal N(\operatorname{vec}W,I_m\otimes\Sigma_U),\\
 T_{lj}:=\tr(\Pi_l^0\Pi_j(U)),\quad T_{lj}\ge0,\quad
 \sum_jT_{lj}=q_l^0,\quad\sum_lT_{lj}=q_j.
\end{gather}
Then the exact orientation-dependent cost is
\begin{align}
 \operatorname{KL}(Q_U\|P_V)
 &=\tfrac12\tr(W\Lambda_0^{-1}W^\top)
 +\tfrac m2\left[\sum_{l,j}\frac{\sigma_j^2}{\tau_l^2}T_{lj}-d
 +\sum_lq_l^0\log\tau_l^2-\sum_jq_j\log\sigma_j^2\right].
\end{align}
A data-dependent $V$ is admissible if it is learned from an independent pilot
sample $S_0$ and the bound is evaluated on a disjoint iid sample $S_1$ of size $n_1$:
\begin{gather}
 S_0\perp S_1,\quad V=V(S_0),\\
 \mathcal E(S_0):=\left\{\begin{aligned}
 &\forall Q\ll P_{V(S_0)}:\\
 &\operatorname{kl}(\widehat L_{S_1}(Q)\|L(Q))
 \le\frac{\operatorname{KL}(Q\|P_{V(S_0)})+\log(\Xi_{n_1}/\delta)}{n_1}
 \end{aligned}\right\},\\
 \PP_{S_1\mid S_0}(\mathcal E(S_0))\ge1-\delta.
\end{gather}
The posterior may use both samples; the empirical term in this bound uses $S_1$.
\end{proposition}
\begin{proof}
\begin{align}
 \tr(\Lambda_0^{-1}\Sigma_U)
 &=\sum_{l,j}\tau_l^{-2}\sigma_j^2\tr(\Pi_l^0\Pi_j),\\
 T_{lj}&=\tr(\Pi_j\Pi_l^0\Pi_j)=\|\Pi_l^0\Pi_j\|_F^2\ge0,\\
 \log\det\Lambda_0-\log\det\Sigma_U
 &=\sum_lq_l^0\log\tau_l^2-\sum_jq_j\log\sigma_j^2.
\end{align}
Substitution in the Gaussian density-ratio identity proves the first assertion.
Conditioning on $S_0$ makes the prior independent of $S_1$; apply
Theorem~\ref{thm:rigid-pac} conditionally, then integrate over $S_0$.
\end{proof}

